% Generated by code/build_component_catalogue.py; edit the generator, not the ID records.
\section{全成分の定義と一般族の一対一カタログ}\label{sec:component-catalogue}
本節の各IDは一つの閉既約像を定義する。表の次元は射影化後、$\mathrm{PGL}_5$商前である。
以下の各source $S_i$と写像$\rho_i$から
\[
 C_i=\overline{\mathrm{PGL}_5\cdot\rho_i(S_i)}_{\mathrm{red}}
\]
を作る。ここでは全軌道の代表元を一つ選ぶのでなく、各成分に稠密な既約族を指定する。
同じ構成式の異なる正規形・枝を同じsourceに併合しない。
各行の「源」は指定命題中の全パラメータ表示とその行の枝である。
「全磁気源」では任意二次式の六係数を全て残す。
座標frame、定数scale、変数平行移動を指定命題が許す範囲で戻す。
非零テンソルの開を取り、rank四行ではさらに$\pi^2\ne0$の開を取る。
これらは有限個の係数の非消滅条件であり、源の既約性と像閉包を変えない。
逆の全点被覆はこの開の表示だけから推論せず、各入口の全支持定理と命題\ref{prop:global-exact147}で証明する。

\subsection{51族を与える有限有理演算の定義}\label{source:circuits51}
51族では、添付の有限有理演算列を数式の圧縮表示として用いる。
これは定義イデアルの未実行の消去を指示するものではない。
添字は零始まりとする。各IDのJSONにある\texttt{nodes}を順に評価し、
\[
 q(\texttt{a/b})=a/b,\quad v(j)=t_j,\quad +(j_1,\ldots,j_k)=\sum z_{j_\ell},\quad
 *(j_1,\ldots,j_k)=\prod z_{j_\ell},\quad \operatorname{inv}(j)=z_j^{-1}
\]
で$z_0,z_1,\ldots$を定める。参照節点は常に前の節点である。
各\texttt{inv}の引数が非零である開を$U_i\subset\mathbb A^{n_i}$とし、
JSONの\texttt{output\_u1\_through\_u61}に指定した61節点を$u_1,\ldots,u_{61}$とする。
演算列に別途記載したperiod Pfaffian及び係数分母の非零条件も課す。
DPMのみ定数$r=\sqrt3\in\CC$を一つ固定して用いる。他方の根を別成分として数えない。
全ての源は空でないaffine空間の開であり既約である。

61値を全150係数に戻す規則も固定する。
対の順は$(01,02,03,04,12,13,14,23,24,34)$、単項式は
$x_ax_b$ ($0\le a\le b\le4$)の辞書式順とし、対を外側に置く。
初期値を$x_0^2(\partial_1\wedge\partial_2+\partial_3\wedge\partial_4)$の係数とする。
\path{legacy/v02/data/symplectic_slice_certificate.json}の
\texttt{raw\_free\_indices}の第$(j+1)$要素が指す係数位置に
$u_{j+1}$ ($0\le j<61$)を置き、
\path{legacy/v02/data/symplectic_slice_equations.json}の
\texttt{eliminated\_raw\_variables}の各$(k,p_k)$について係数$k$を$p_k(u)$へ置換する。
多項式の符号化$(J,a,b)$は$(a/b)\prod_{j\in J}u_{j+1}$を意味する。
他の係数は初期値のままである。この規則を$\mathcal R(u)$と書く。
\[
 \rho_i(t)=[\mathcal R(u(t))]\in\mathbb P H^0(\wedge^2T).
\]
実装は\path{code/s61_model.py}の\texttt{S61.reconstruct}と同じ規則であり、
座標の定義は実装の選択に依存しない。

T01--T38はperiod、装飾辺の平滑化係数及び$\mathrm{Sp}_4$の大cell十係数を入力とする。
全periodは装飾辺の次数とforced-zero条件の線形核で与えられ、
全有限Maurer--Cartan補正を含めた演算列を用いる。
T39,T40は十六一次係数からの既存の二次graph表示である。
トーリック境界のbiresidue行列だけを一般の平滑化元と取り違えない。
MPSのTheorem1.4、Theorem6.16、Corollary6.23とExample6.24の局所Torelliによる
成分性・相違を用い、Table2との全装飾図式対応を固定している\cite{mps}。
この一般族のPoisson性は全sourceのJacobi構成と実射影変換による。
源の稠密性は、一般rank四点を選び値を標準化した後の全symplectic frameを
大cellで尽くし、perspective四方向を一意に除くことから従う。
追加十一族では一般点四係数と全体scaleも残し、同じ論法を使う。
その成分性・相違は本文の全周囲横断切片証明による。
根拠と演算列のhashは
\path{code/data/known51_rank4_s61_coordinate_catalogue_20261103_v62.json}、
一般恒等式と正規化の証明は
\path{research/notes/MPS_S61_CATALOGUE_20261103_v62.md}及び
\path{research/notes/KNOWN51_RANK4_COORDINATE_MAPS_20261103_v62.md}にある。
正確な標本評価は実装検査であり、この全パラメータの証明の代用ではない。

\subsection{NEW25の三変数多項式源}\label{source:new25}
$(x,y,z,w)$をaffine座標とし、$n_0,n_1,n_2$を任意とする。
六係数を次のように置き、$h$で次数二に同次化する。
\[
\begin{aligned}
 P_{xy}&=n_1/36-256w/3,&P_{xw}&=8/9,\\
 P_{xz}&=n_1n_2y/48+n_1x/48-64n_2wy-64wx,\\
 P_{yz}&=n_0n_2y^2/2+n_0xy-4n_0z/3-n_1^2w/1536\\
 &\quad+n_1w^2-n_1x^2/64+64wy-2/9,\\
 P_{yw}&=4x/3,&P_{zw}&=n_2xy+x^2+2y/3.
\end{aligned}
\]
これは全$\mathbb A^3$上のPoisson族であり、$n_0^2n_2=0$という補助条件を課さない。
以下の全150係数へのEuler射影は射影sectionを変えない。
一般因子は$3H+Q$、$Q$は滑らかで$H$は接する。
成分性の上界と相違は本文のNEW25の八初期式の証明による。

\subsection{最後の五成分の明示源}\label{source:last-five-explicit}
最後の五族について、参照先の全支持証明とは別に構成式をここで固定する。
まず$(h,x,y,z,w)$でhをCasimirとし、未表示括弧を零とする。
Ga9 linear-wの源は一点$P_G$で、
\[
\begin{aligned}
 P_G^{xy}&=4h^2-y^2,&P_G^{xz}&=-hx+2yz+4w^2,\\
 P_G^{xw}&=4y(4h-w),&P_G^{yz}&=hy,\\
 P_G^{zw}&=4(y^2-hw).
\end{aligned}
\]
次の三源では$R=\sum_{i+j+k=2}r_{ijk}h^ix^jy^k$の全六係数を任意とする。
順にnode変曲接線・変曲点を通る通常接線、cusp変曲接線・変曲点を通る割線、
cusp接線・cuspを通る割線のsourceである。
\[
\begin{aligned}
 P_1^{xz}&=2x(x-h),&P_1^{xw}&=-2xy,\\
 P_1^{yz}&=y(3x-h),&P_1^{yw}&=2x(h-x)-3y^2,\\
 P_1^{zw}&=-(x-h)w+R;
\end{aligned}
\]
\[
\begin{aligned}
 P_2^{xz}&=2x(x+h),&P_2^{xw}&=2xy,\\
 P_2^{yz}&=3xy,&P_2^{yw}&=3(y^2-(h+x)^2),\\
 P_2^{zw}&=-xw+R;
\end{aligned}
\]
\[
\begin{aligned}
 P_3^{xz}&=-hx,&P_3^{xw}&=3x(h-x),\\
 P_3^{yz}&=3(h-x)^2-3hy,&P_3^{yw}&=-9xy,\\
 P_3^{zw}&=3xz+3R.
\end{aligned}
\]
最後に$(h,t,x,y,r)$でhをCasimirとし、一点$P_D$を
\[
\begin{aligned}
 P_D^{tx}&=h^2+hx-2xy,&P_D^{ty}&=y(4y-h),\\
 P_D^{tr}&=(h-6y)r+hx,&P_D^{xr}&=-x^2,\\
 P_D^{yr}&=h^2+2xy
\end{aligned}
\]
とする。$P_G,P_D$には$S=\operatorname{Spec}\CC$、三磁気族には$S=\mathbb A^6$を用い、
いずれも全$\mathrm{PGL}_5$飽和閉包を取る。
全sourceでJacobiが成立し実Pfaffianが非零である。
その全源の認識、次元、成分性・相違は命題\ref{prop:five-last-full-components}と参照する全支持証明による。
三磁気族を一つの標本点の軌道で置き換えていない。

\subsection{全IDの定義・幾何・証明}
$D_d$は次数$d$の既約超曲面、$H_i$は相異なる超平面を意味する。
同じ次数が複数並ぶときも相異なる因子であり、指数が重複度である。
rank二行ではactual退化五次は恒等零なので因子型を付けず、foliationとその零因子を記す。
全IDの支持への参照は命題\ref{prop:global-exact147}、
極大性の共通判定は補題\ref{lem:transverse-support}と
第\ref{sec:maximal-image-synthesis}節である。

\paragraph{族1}\label{catalogue:T01}
\path{T01}。
射影次元$25$、一般階数$4$。
源と写像は第\ref{source:circuits51}節の有理演算源$U_i$、$\rho_i$。演算列は\path{reports/MPS_S61_FULL_CIRCUITS_20261103_v62/T01.json.gz}。
入力は\path{p0,p1,p2,p3,p4,p5,g0,g1,g2,g3,g4,g5,g6,g7,g8,g9}。全inverse節点と記載した分母の非零開を取る。
一般実退化因子は$H_{1}+H_{2}+H_{3}+H_{4}+H_{5}$。
構成は第\ref{source:circuits51}節、成分性は第\ref{source:circuits51}節、相違は第\ref{source:circuits51}節及びその証明を参照する。

\paragraph{族2}\label{catalogue:T02}
\path{T02}。
射影次元$24$、一般階数$4$。
源と写像は第\ref{source:circuits51}節の有理演算源$U_i$、$\rho_i$。演算列は\path{reports/MPS_S61_FULL_CIRCUITS_20261103_v62/T02.json.gz}。
入力は\path{p0,p1,p2,p3,t0,g0,g1,g2,g3,g4,g5,g6,g7,g8,g9}。全inverse節点と記載した分母の非零開を取る。
一般実退化因子は$D_{2,1}+H_{2}+H_{3}+H_{4}$。
構成は第\ref{source:circuits51}節、成分性は第\ref{source:circuits51}節、相違は第\ref{source:circuits51}節及びその証明を参照する。

\paragraph{族3}\label{catalogue:T03}
\path{T03}。
射影次元$24$、一般階数$4$。
源と写像は第\ref{source:circuits51}節の有理演算源$U_i$、$\rho_i$。演算列は\path{reports/MPS_S61_FULL_CIRCUITS_20261103_v62/T03.json.gz}。
入力は\path{p0,p1,p2,p3,t0,g0,g1,g2,g3,g4,g5,g6,g7,g8,g9}。全inverse節点と記載した分母の非零開を取る。
一般実退化因子は$D_{2,1}+H_{2}+H_{3}+H_{4}$。
構成は第\ref{source:circuits51}節、成分性は第\ref{source:circuits51}節、相違は第\ref{source:circuits51}節及びその証明を参照する。

\paragraph{族4}\label{catalogue:T04}
\path{T04}。
射影次元$23$、一般階数$4$。
源と写像は第\ref{source:circuits51}節の有理演算源$U_i$、$\rho_i$。演算列は\path{reports/MPS_S61_FULL_CIRCUITS_20261103_v62/T04.json.gz}。
入力は\path{p0,p1,t0,t1,g0,g1,g2,g3,g4,g5,g6,g7,g8,g9}。全inverse節点と記載した分母の非零開を取る。
一般実退化因子は$D_{3,1}+H_{2}+H_{3}$。
構成は第\ref{source:circuits51}節、成分性は第\ref{source:circuits51}節、相違は第\ref{source:circuits51}節及びその証明を参照する。

\paragraph{族5}\label{catalogue:T05}
\path{T05}。
射影次元$23$、一般階数$4$。
源と写像は第\ref{source:circuits51}節の有理演算源$U_i$、$\rho_i$。演算列は\path{reports/MPS_S61_FULL_CIRCUITS_20261103_v62/T05.json.gz}。
入力は\path{p0,p1,t0,t1,g0,g1,g2,g3,g4,g5,g6,g7,g8,g9}。全inverse節点と記載した分母の非零開を取る。
一般実退化因子は$D_{3,1}+H_{2}+H_{3}$。
構成は第\ref{source:circuits51}節、成分性は第\ref{source:circuits51}節、相違は第\ref{source:circuits51}節及びその証明を参照する。

\paragraph{族6}\label{catalogue:T06}
\path{T06}。
射影次元$23$、一般階数$4$。
源と写像は第\ref{source:circuits51}節の有理演算源$U_i$、$\rho_i$。演算列は\path{reports/MPS_S61_FULL_CIRCUITS_20261103_v62/T06.json.gz}。
入力は\path{p0,p1,t0,t1,g0,g1,g2,g3,g4,g5,g6,g7,g8,g9}。全inverse節点と記載した分母の非零開を取る。
一般実退化因子は$D_{3,1}+H_{2}+H_{3}$。
構成は第\ref{source:circuits51}節、成分性は第\ref{source:circuits51}節、相違は第\ref{source:circuits51}節及びその証明を参照する。

\paragraph{族7}\label{catalogue:T07}
\path{T07}。
射影次元$23$、一般階数$4$。
源と写像は第\ref{source:circuits51}節の有理演算源$U_i$、$\rho_i$。演算列は\path{reports/MPS_S61_FULL_CIRCUITS_20261103_v62/T07.json.gz}。
入力は\path{p0,p1,t0,t1,g0,g1,g2,g3,g4,g5,g6,g7,g8,g9}。全inverse節点と記載した分母の非零開を取る。
一般実退化因子は$D_{3,1}+H_{2}+H_{3}$。
構成は第\ref{source:circuits51}節、成分性は第\ref{source:circuits51}節、相違は第\ref{source:circuits51}節及びその証明を参照する。

\paragraph{族8}\label{catalogue:T08}
\path{T08}。
射影次元$23$、一般階数$4$。
源と写像は第\ref{source:circuits51}節の有理演算源$U_i$、$\rho_i$。演算列は\path{reports/MPS_S61_FULL_CIRCUITS_20261103_v62/T08.json.gz}。
入力は\path{p0,p1,t0,t1,g0,g1,g2,g3,g4,g5,g6,g7,g8,g9}。全inverse節点と記載した分母の非零開を取る。
一般実退化因子は$D_{3,1}+H_{2}+H_{3}$。
構成は第\ref{source:circuits51}節、成分性は第\ref{source:circuits51}節、相違は第\ref{source:circuits51}節及びその証明を参照する。

\paragraph{族9}\label{catalogue:T09}
\path{T09}。
射影次元$23$、一般階数$4$。
源と写像は第\ref{source:circuits51}節の有理演算源$U_i$、$\rho_i$。演算列は\path{reports/MPS_S61_FULL_CIRCUITS_20261103_v62/T09.json.gz}。
入力は\path{p0,p1,t0,t1,g0,g1,g2,g3,g4,g5,g6,g7,g8,g9}。全inverse節点と記載した分母の非零開を取る。
一般実退化因子は$D_{3,1}+H_{2}+H_{3}$。
構成は第\ref{source:circuits51}節、成分性は第\ref{source:circuits51}節、相違は第\ref{source:circuits51}節及びその証明を参照する。

\paragraph{族10}\label{catalogue:T10}
\path{T10}。
射影次元$23$、一般階数$4$。
源と写像は第\ref{source:circuits51}節の有理演算源$U_i$、$\rho_i$。演算列は\path{reports/MPS_S61_FULL_CIRCUITS_20261103_v62/T10.json.gz}。
入力は\path{p0,p1,t0,t1,g0,g1,g2,g3,g4,g5,g6,g7,g8,g9}。全inverse節点と記載した分母の非零開を取る。
一般実退化因子は$D_{3,1}+H_{2}+H_{3}$。
構成は第\ref{source:circuits51}節、成分性は第\ref{source:circuits51}節、相違は第\ref{source:circuits51}節及びその証明を参照する。

\paragraph{族11}\label{catalogue:T11}
\path{T11}。
射影次元$23$、一般階数$4$。
源と写像は第\ref{source:circuits51}節の有理演算源$U_i$、$\rho_i$。演算列は\path{reports/MPS_S61_FULL_CIRCUITS_20261103_v62/T11.json.gz}。
入力は\path{p0,p1,t0,t1,g0,g1,g2,g3,g4,g5,g6,g7,g8,g9}。全inverse節点と記載した分母の非零開を取る。
一般実退化因子は$D_{3,1}+H_{2}+H_{3}$。
構成は第\ref{source:circuits51}節、成分性は第\ref{source:circuits51}節、相違は第\ref{source:circuits51}節及びその証明を参照する。

\paragraph{族12}\label{catalogue:T12}
\path{T12}。
射影次元$23$、一般階数$4$。
源と写像は第\ref{source:circuits51}節の有理演算源$U_i$、$\rho_i$。演算列は\path{reports/MPS_S61_FULL_CIRCUITS_20261103_v62/T12.json.gz}。
入力は\path{p0,p1,t0,t1,g0,g1,g2,g3,g4,g5,g6,g7,g8,g9}。全inverse節点と記載した分母の非零開を取る。
一般実退化因子は$D_{3,1}+H_{2}+H_{3}$。
構成は第\ref{source:circuits51}節、成分性は第\ref{source:circuits51}節、相違は第\ref{source:circuits51}節及びその証明を参照する。

\paragraph{族13}\label{catalogue:T13}
\path{T13}。
射影次元$23$、一般階数$4$。
源と写像は第\ref{source:circuits51}節の有理演算源$U_i$、$\rho_i$。演算列は\path{reports/MPS_S61_FULL_CIRCUITS_20261103_v62/T13.json.gz}。
入力は\path{p0,p1,t0,t1,g0,g1,g2,g3,g4,g5,g6,g7,g8,g9}。全inverse節点と記載した分母の非零開を取る。
一般実退化因子は$D_{3,1}+H_{2}+H_{3}$。
構成は第\ref{source:circuits51}節、成分性は第\ref{source:circuits51}節、相違は第\ref{source:circuits51}節及びその証明を参照する。

\paragraph{族14}\label{catalogue:T14}
\path{T14}。
射影次元$23$、一般階数$4$。
源と写像は第\ref{source:circuits51}節の有理演算源$U_i$、$\rho_i$。演算列は\path{reports/MPS_S61_FULL_CIRCUITS_20261103_v62/T14.json.gz}。
入力は\path{p0,p1,t0,t1,g0,g1,g2,g3,g4,g5,g6,g7,g8,g9}。全inverse節点と記載した分母の非零開を取る。
一般実退化因子は$D_{3,1}+H_{2}+H_{3}$。
構成は第\ref{source:circuits51}節、成分性は第\ref{source:circuits51}節、相違は第\ref{source:circuits51}節及びその証明を参照する。

\paragraph{族15}\label{catalogue:T15}
\path{T15}。
射影次元$23$、一般階数$4$。
源と写像は第\ref{source:circuits51}節の有理演算源$U_i$、$\rho_i$。演算列は\path{reports/MPS_S61_FULL_CIRCUITS_20261103_v62/T15.json.gz}。
入力は\path{p0,p1,t0,t1,g0,g1,g2,g3,g4,g5,g6,g7,g8,g9}。全inverse節点と記載した分母の非零開を取る。
一般実退化因子は$D_{3,1}+H_{2}+H_{3}$。
構成は第\ref{source:circuits51}節、成分性は第\ref{source:circuits51}節、相違は第\ref{source:circuits51}節及びその証明を参照する。

\paragraph{族16}\label{catalogue:T16}
\path{T16}。
射影次元$24$、一般階数$4$。
源と写像は第\ref{source:circuits51}節の有理演算源$U_i$、$\rho_i$。演算列は\path{reports/MPS_S61_FULL_CIRCUITS_20261103_v62/T16.json.gz}。
入力は\path{p0,p1,p2,t0,t1,g0,g1,g2,g3,g4,g5,g6,g7,g8,g9}。全inverse節点と記載した分母の非零開を取る。
一般実退化因子は$D_{2,1}+D_{2,2}+H_{3}$。
構成は第\ref{source:circuits51}節、成分性は第\ref{source:circuits51}節、相違は第\ref{source:circuits51}節及びその証明を参照する。

\paragraph{族17}\label{catalogue:T17}
\path{T17}。
射影次元$24$、一般階数$4$。
源と写像は第\ref{source:circuits51}節の有理演算源$U_i$、$\rho_i$。演算列は\path{reports/MPS_S61_FULL_CIRCUITS_20261103_v62/T17.json.gz}。
入力は\path{p0,p1,p2,t0,t1,g0,g1,g2,g3,g4,g5,g6,g7,g8,g9}。全inverse節点と記載した分母の非零開を取る。
一般実退化因子は$D_{2,1}+D_{2,2}+H_{3}$。
構成は第\ref{source:circuits51}節、成分性は第\ref{source:circuits51}節、相違は第\ref{source:circuits51}節及びその証明を参照する。

\paragraph{族18}\label{catalogue:T18}
\path{T18}。
射影次元$23$、一般階数$4$。
源と写像は第\ref{source:circuits51}節の有理演算源$U_i$、$\rho_i$。演算列は\path{reports/MPS_S61_FULL_CIRCUITS_20261103_v62/T18.json.gz}。
入力は\path{p0,p1,t0,t1,g0,g1,g2,g3,g4,g5,g6,g7,g8,g9}。全inverse節点と記載した分母の非零開を取る。
一般実退化因子は$D_{2,1}+D_{2,2}+H_{3}$。
構成は第\ref{source:circuits51}節、成分性は第\ref{source:circuits51}節、相違は第\ref{source:circuits51}節及びその証明を参照する。

\paragraph{族19}\label{catalogue:T19}
\path{T19}。
射影次元$23$、一般階数$4$。
源と写像は第\ref{source:circuits51}節の有理演算源$U_i$、$\rho_i$。演算列は\path{reports/MPS_S61_FULL_CIRCUITS_20261103_v62/T19.json.gz}。
入力は\path{p0,p1,t0,t1,g0,g1,g2,g3,g4,g5,g6,g7,g8,g9}。全inverse節点と記載した分母の非零開を取る。
一般実退化因子は$D_{2,1}+D_{2,2}+H_{3}$。
構成は第\ref{source:circuits51}節、成分性は第\ref{source:circuits51}節、相違は第\ref{source:circuits51}節及びその証明を参照する。

\paragraph{族20}\label{catalogue:T20}
\path{T20}。
射影次元$23$、一般階数$4$。
源と写像は第\ref{source:circuits51}節の有理演算源$U_i$、$\rho_i$。演算列は\path{reports/MPS_S61_FULL_CIRCUITS_20261103_v62/T20.json.gz}。
入力は\path{p0,p1,t0,t1,g0,g1,g2,g3,g4,g5,g6,g7,g8,g9}。全inverse節点と記載した分母の非零開を取る。
一般実退化因子は$D_{2,1}+D_{2,2}+H_{3}$。
構成は第\ref{source:circuits51}節、成分性は第\ref{source:circuits51}節、相違は第\ref{source:circuits51}節及びその証明を参照する。

\paragraph{族21}\label{catalogue:T21}
\path{T21}。
射影次元$24$、一般階数$4$。
源と写像は第\ref{source:circuits51}節の有理演算源$U_i$、$\rho_i$。演算列は\path{reports/MPS_S61_FULL_CIRCUITS_20261103_v62/T21.json.gz}。
入力は\path{p0,p1,t0,t1,t2,g0,g1,g2,g3,g4,g5,g6,g7,g8,g9}。全inverse節点と記載した分母の非零開を取る。
一般実退化因子は$D_{3,1}+H_{2}+H_{3}$。
構成は第\ref{source:circuits51}節、成分性は第\ref{source:circuits51}節、相違は第\ref{source:circuits51}節及びその証明を参照する。

\paragraph{族22}\label{catalogue:T22}
\path{T22}。
射影次元$23$、一般階数$4$。
源と写像は第\ref{source:circuits51}節の有理演算源$U_i$、$\rho_i$。演算列は\path{reports/MPS_S61_FULL_CIRCUITS_20261103_v62/T22.json.gz}。
入力は\path{p0,t0,t1,t2,g0,g1,g2,g3,g4,g5,g6,g7,g8,g9}。全inverse節点と記載した分母の非零開を取る。
一般実退化因子は$D_{4,1}+H_{2}$。
構成は第\ref{source:circuits51}節、成分性は第\ref{source:circuits51}節、相違は第\ref{source:circuits51}節及びその証明を参照する。

\paragraph{族23}\label{catalogue:T23}
\path{T23}。
射影次元$23$、一般階数$4$。
源と写像は第\ref{source:circuits51}節の有理演算源$U_i$、$\rho_i$。演算列は\path{reports/MPS_S61_FULL_CIRCUITS_20261103_v62/T23.json.gz}。
入力は\path{p0,t0,t1,t2,g0,g1,g2,g3,g4,g5,g6,g7,g8,g9}。全inverse節点と記載した分母の非零開を取る。
一般実退化因子は$D_{4,1}+H_{2}$。
構成は第\ref{source:circuits51}節、成分性は第\ref{source:circuits51}節、相違は第\ref{source:circuits51}節及びその証明を参照する。

\paragraph{族24}\label{catalogue:T24}
\path{T24}。
射影次元$23$、一般階数$4$。
源と写像は第\ref{source:circuits51}節の有理演算源$U_i$、$\rho_i$。演算列は\path{reports/MPS_S61_FULL_CIRCUITS_20261103_v62/T24.json.gz}。
入力は\path{p0,t0,t1,t2,g0,g1,g2,g3,g4,g5,g6,g7,g8,g9}。全inverse節点と記載した分母の非零開を取る。
一般実退化因子は$D_{4,1}+H_{2}$。
構成は第\ref{source:circuits51}節、成分性は第\ref{source:circuits51}節、相違は第\ref{source:circuits51}節及びその証明を参照する。

\paragraph{族25}\label{catalogue:T25}
\path{T25}。
射影次元$23$、一般階数$4$。
源と写像は第\ref{source:circuits51}節の有理演算源$U_i$、$\rho_i$。演算列は\path{reports/MPS_S61_FULL_CIRCUITS_20261103_v62/T25.json.gz}。
入力は\path{p0,t0,t1,t2,g0,g1,g2,g3,g4,g5,g6,g7,g8,g9}。全inverse節点と記載した分母の非零開を取る。
一般実退化因子は$D_{4,1}+H_{2}$。
構成は第\ref{source:circuits51}節、成分性は第\ref{source:circuits51}節、相違は第\ref{source:circuits51}節及びその証明を参照する。

\paragraph{族26}\label{catalogue:T26}
\path{T26}。
射影次元$23$、一般階数$4$。
源と写像は第\ref{source:circuits51}節の有理演算源$U_i$、$\rho_i$。演算列は\path{reports/MPS_S61_FULL_CIRCUITS_20261103_v62/T26.json.gz}。
入力は\path{p0,t0,t1,t2,g0,g1,g2,g3,g4,g5,g6,g7,g8,g9}。全inverse節点と記載した分母の非零開を取る。
一般実退化因子は$D_{4,1}+H_{2}$。
構成は第\ref{source:circuits51}節、成分性は第\ref{source:circuits51}節、相違は第\ref{source:circuits51}節及びその証明を参照する。

\paragraph{族27}\label{catalogue:T27}
\path{T27}。
射影次元$23$、一般階数$4$。
源と写像は第\ref{source:circuits51}節の有理演算源$U_i$、$\rho_i$。演算列は\path{reports/MPS_S61_FULL_CIRCUITS_20261103_v62/T27.json.gz}。
入力は\path{p0,t0,t1,t2,g0,g1,g2,g3,g4,g5,g6,g7,g8,g9}。全inverse節点と記載した分母の非零開を取る。
一般実退化因子は$D_{4,1}+H_{2}$。
構成は第\ref{source:circuits51}節、成分性は第\ref{source:circuits51}節、相違は第\ref{source:circuits51}節及びその証明を参照する。

\paragraph{族28}\label{catalogue:T28}
\path{T28}。
射影次元$23$、一般階数$4$。
源と写像は第\ref{source:circuits51}節の有理演算源$U_i$、$\rho_i$。演算列は\path{reports/MPS_S61_FULL_CIRCUITS_20261103_v62/T28.json.gz}。
入力は\path{p0,t0,t1,t2,g0,g1,g2,g3,g4,g5,g6,g7,g8,g9}。全inverse節点と記載した分母の非零開を取る。
一般実退化因子は$D_{4,1}+H_{2}$。
構成は第\ref{source:circuits51}節、成分性は第\ref{source:circuits51}節、相違は第\ref{source:circuits51}節及びその証明を参照する。

\paragraph{族29}\label{catalogue:T29}
\path{T29}。
射影次元$23$、一般階数$4$。
源と写像は第\ref{source:circuits51}節の有理演算源$U_i$、$\rho_i$。演算列は\path{reports/MPS_S61_FULL_CIRCUITS_20261103_v62/T29.json.gz}。
入力は\path{p0,t0,t1,t2,g0,g1,g2,g3,g4,g5,g6,g7,g8,g9}。全inverse節点と記載した分母の非零開を取る。
一般実退化因子は$D_{4,1}+H_{2}$。
構成は第\ref{source:circuits51}節、成分性は第\ref{source:circuits51}節、相違は第\ref{source:circuits51}節及びその証明を参照する。

\paragraph{族30}\label{catalogue:T30}
\path{T30}。
射影次元$23$、一般階数$4$。
源と写像は第\ref{source:circuits51}節の有理演算源$U_i$、$\rho_i$。演算列は\path{reports/MPS_S61_FULL_CIRCUITS_20261103_v62/T30.json.gz}。
入力は\path{p0,t0,t1,t2,g0,g1,g2,g3,g4,g5,g6,g7,g8,g9}。全inverse節点と記載した分母の非零開を取る。
一般実退化因子は$D_{4,1}+H_{2}$。
構成は第\ref{source:circuits51}節、成分性は第\ref{source:circuits51}節、相違は第\ref{source:circuits51}節及びその証明を参照する。

\paragraph{族31}\label{catalogue:T31}
\path{T31}。
射影次元$23$、一般階数$4$。
源と写像は第\ref{source:circuits51}節の有理演算源$U_i$、$\rho_i$。演算列は\path{reports/MPS_S61_FULL_CIRCUITS_20261103_v62/T31.json.gz}。
入力は\path{p0,t0,t1,t2,g0,g1,g2,g3,g4,g5,g6,g7,g8,g9}。全inverse節点と記載した分母の非零開を取る。
一般実退化因子は$D_{4,1}+H_{2}$。
構成は第\ref{source:circuits51}節、成分性は第\ref{source:circuits51}節、相違は第\ref{source:circuits51}節及びその証明を参照する。

\paragraph{族32}\label{catalogue:T32}
\path{T32}。
射影次元$23$、一般階数$4$。
源と写像は第\ref{source:circuits51}節の有理演算源$U_i$、$\rho_i$。演算列は\path{reports/MPS_S61_FULL_CIRCUITS_20261103_v62/T32.json.gz}。
入力は\path{p0,t0,t1,t2,g0,g1,g2,g3,g4,g5,g6,g7,g8,g9}。全inverse節点と記載した分母の非零開を取る。
一般実退化因子は$D_{3,1}+D_{2,2}$。
構成は第\ref{source:circuits51}節、成分性は第\ref{source:circuits51}節、相違は第\ref{source:circuits51}節及びその証明を参照する。

\paragraph{族33}\label{catalogue:T33}
\path{T33}。
射影次元$23$、一般階数$4$。
源と写像は第\ref{source:circuits51}節の有理演算源$U_i$、$\rho_i$。演算列は\path{reports/MPS_S61_FULL_CIRCUITS_20261103_v62/T33.json.gz}。
入力は\path{p0,t0,t1,t2,g0,g1,g2,g3,g4,g5,g6,g7,g8,g9}。全inverse節点と記載した分母の非零開を取る。
一般実退化因子は$D_{3,1}+D_{2,2}$。
構成は第\ref{source:circuits51}節、成分性は第\ref{source:circuits51}節、相違は第\ref{source:circuits51}節及びその証明を参照する。

\paragraph{族34}\label{catalogue:T34}
\path{T34}。
射影次元$23$、一般階数$4$。
源と写像は第\ref{source:circuits51}節の有理演算源$U_i$、$\rho_i$。演算列は\path{reports/MPS_S61_FULL_CIRCUITS_20261103_v62/T34.json.gz}。
入力は\path{p0,t0,t1,t2,g0,g1,g2,g3,g4,g5,g6,g7,g8,g9}。全inverse節点と記載した分母の非零開を取る。
一般実退化因子は$D_{3,1}+D_{2,2}$。
構成は第\ref{source:circuits51}節、成分性は第\ref{source:circuits51}節、相違は第\ref{source:circuits51}節及びその証明を参照する。

\paragraph{族35}\label{catalogue:T35}
\path{T35}。
射影次元$23$、一般階数$4$。
源と写像は第\ref{source:circuits51}節の有理演算源$U_i$、$\rho_i$。演算列は\path{reports/MPS_S61_FULL_CIRCUITS_20261103_v62/T35.json.gz}。
入力は\path{p0,t0,t1,t2,g0,g1,g2,g3,g4,g5,g6,g7,g8,g9}。全inverse節点と記載した分母の非零開を取る。
一般実退化因子は$D_{3,1}+D_{2,2}$。
構成は第\ref{source:circuits51}節、成分性は第\ref{source:circuits51}節、相違は第\ref{source:circuits51}節及びその証明を参照する。

\paragraph{族36}\label{catalogue:T36}
\path{T36}。
射影次元$24$、一般階数$4$。
源と写像は第\ref{source:circuits51}節の有理演算源$U_i$、$\rho_i$。演算列は\path{reports/MPS_S61_FULL_CIRCUITS_20261103_v62/T36.json.gz}。
入力は\path{p0,t0,t1,t2,t3,g0,g1,g2,g3,g4,g5,g6,g7,g8,g9}。全inverse節点と記載した分母の非零開を取る。
一般実退化因子は$D_{5,1}$。
構成は第\ref{source:circuits51}節、成分性は第\ref{source:circuits51}節、相違は第\ref{source:circuits51}節及びその証明を参照する。

\paragraph{族37}\label{catalogue:T37}
\path{T37}。
射影次元$24$、一般階数$4$。
源と写像は第\ref{source:circuits51}節の有理演算源$U_i$、$\rho_i$。演算列は\path{reports/MPS_S61_FULL_CIRCUITS_20261103_v62/T37.json.gz}。
入力は\path{p0,t0,t1,t2,t3,g0,g1,g2,g3,g4,g5,g6,g7,g8,g9}。全inverse節点と記載した分母の非零開を取る。
一般実退化因子は$D_{5,1}$。
構成は第\ref{source:circuits51}節、成分性は第\ref{source:circuits51}節、相違は第\ref{source:circuits51}節及びその証明を参照する。

\paragraph{族38}\label{catalogue:T38}
\path{T38}。
射影次元$24$、一般階数$4$。
源と写像は第\ref{source:circuits51}節の有理演算源$U_i$、$\rho_i$。演算列は\path{reports/MPS_S61_FULL_CIRCUITS_20261103_v62/T38.json.gz}。
入力は\path{p0,t0,t1,t2,t3,g0,g1,g2,g3,g4,g5,g6,g7,g8,g9}。全inverse節点と記載した分母の非零開を取る。
一般実退化因子は$D_{5,1}$。
構成は第\ref{source:circuits51}節、成分性は第\ref{source:circuits51}節、相違は第\ref{source:circuits51}節及びその証明を参照する。

\paragraph{族39}\label{catalogue:T39}
\path{T39}。
射影次元$25$、一般階数$4$。
源と写像は第\ref{source:circuits51}節の有理演算源$U_i$、$\rho_i$。演算列は\path{reports/MPS_ELLIPTIC_GRAPH_CIRCUITS_20261103_v62/T39.json.gz}。
入力は\path{u4,u5,u12,u13,u14,u25,u26,u27,u34,u35,u36,u37,u48,u49,u50,u51}。全inverse節点と記載した分母の非零開を取る。
一般実退化因子は$D_{5,1}$。
構成は第\ref{source:circuits51}節、成分性は第\ref{source:circuits51}節、相違は第\ref{source:circuits51}節及びその証明を参照する。

\paragraph{族40}\label{catalogue:T40}
\path{T40}。
射影次元$25$、一般階数$4$。
源と写像は第\ref{source:circuits51}節の有理演算源$U_i$、$\rho_i$。演算列は\path{reports/MPS_ELLIPTIC_GRAPH_CIRCUITS_20261103_v62/T40.json.gz}。
入力は\path{u4,u5,u12,u13,u14,u25,u26,u27,u34,u35,u36,u37,u48,u49,u50,u51}。全inverse節点と記載した分母の非零開を取る。
一般実退化因子は$D_{5,1}$。
構成は第\ref{source:circuits51}節、成分性は第\ref{source:circuits51}節、相違は第\ref{source:circuits51}節及びその証明を参照する。

\paragraph{族41}\label{catalogue:R2_RAT113}
\path{R2_RAT113}。
射影次元$36$、一般階数$2$。
一般族は全$(L_1,L_2,G_3)$のJacobian源。第\ref{sec:rank-two-reconstruction}節の全sourceを用いる。一般零因子は空、飽和foliationは次数二である。
構成は第\ref{sec:rank-two-reconstruction}節、成分性は第\ref{sec:rank-two-reconstruction}節、相違は第\ref{sec:rank-two-reconstruction}節及びその証明を参照する。

\paragraph{族42}\label{catalogue:R2_RAT122}
\path{R2_RAT122}。
射影次元$28$、一般階数$2$。
一般族は全$(L,Q_1,Q_2)$のJacobian源。第\ref{sec:rank-two-reconstruction}節の全sourceを用いる。一般零因子は空、飽和foliationは次数二である。
構成は第\ref{sec:rank-two-reconstruction}節、成分性は第\ref{sec:rank-two-reconstruction}節、相違は第\ref{sec:rank-two-reconstruction}節及びその証明を参照する。

\paragraph{族43}\label{catalogue:R2_LOG1112}
\path{R2_LOG1112}。
射影次元$28$、一般階数$2$。
一般族は次数$(1,1,1,2)$の全四式と$\bigwedge^2\ker(1,1,1,2)$の留数源。第\ref{sec:rank-two-reconstruction}節の全sourceを用いる。一般零因子は空、飽和foliationは次数二である。
構成は第\ref{sec:rank-two-reconstruction}節、成分性は第\ref{sec:rank-two-reconstruction}節、相違は第\ref{sec:rank-two-reconstruction}節及びその証明を参照する。

\paragraph{族44}\label{catalogue:R2_LPB}
\path{R2_LPB}。
射影次元$39$、一般階数$2$。
一般族は$N=\partial_4$と$x_0,\ldots,x_3$の全二次場$Y$からの$N\wedge Y$。第\ref{sec:rank-two-reconstruction}節の全sourceを用いる。一般零因子は空、飽和foliationは次数二である。
構成は第\ref{sec:rank-two-reconstruction}節、成分性は第\ref{sec:rank-two-reconstruction}節、相違は第\ref{sec:rank-two-reconstruction}節及びその証明を参照する。

\paragraph{族45}\label{catalogue:R2_AFF5}
\path{R2_AFF5}。
射影次元$22$、一般階数$2$。
一般族はJordan分割$(5)$の$N$と全解$[N,H]=N$からの$N\wedge H$。第\ref{sec:rank-two-reconstruction}節の全sourceを用いる。一般零因子は空、飽和foliationは次数二である。
構成は第\ref{sec:rank-two-reconstruction}節、成分性は第\ref{sec:rank-two-reconstruction}節、相違は第\ref{sec:rank-two-reconstruction}節及びその証明を参照する。

\paragraph{族46}\label{catalogue:R2_AFF41}
\path{R2_AFF41}。
射影次元$22$、一般階数$2$。
一般族はJordan分割$(4,1)$の$N$と全解$[N,H]=N$からの$N\wedge H$。第\ref{sec:rank-two-reconstruction}節の全sourceを用いる。一般零因子は空、飽和foliationは次数二である。
構成は第\ref{sec:rank-two-reconstruction}節、成分性は第\ref{sec:rank-two-reconstruction}節、相違は第\ref{sec:rank-two-reconstruction}節及びその証明を参照する。

\paragraph{族47}\label{catalogue:R4_DOUBLE_SMOOTH_QUADRIC}
\path{R4_DOUBLE_SMOOTH_QUADRIC}。
射影次元$22$、一般階数$4$。
源と写像は第\ref{source:circuits51}節の有理演算源$U_i$、$\rho_i$。演算列は\path{reports/ADDITIONAL_RANK4_S61_CIRCUITS_20261103_v62/R4_DOUBLE_SMOOTH_QUADRIC.json.gz}。
入力は\path{scale,base1,base2,base3,base4,g0,g1,g2,g3,g4,g5,g6,g7,g8,g9}。全inverse節点と記載した分母の非零開を取る。
一般実退化因子は$H_{1}+D_{2,2}^{2}$。
description：H+2Q; Q smooth and H tangent to Q。
構成は第\ref{source:circuits51}節、成分性は第\ref{sec:known-component-certificates}節、相違は第\ref{sec:known-component-certificates}節及びその証明を参照する。

\paragraph{族48}\label{catalogue:R4_TRIPLE_HYPERPLANE_QUADRIC}
\path{R4_TRIPLE_HYPERPLANE_QUADRIC}。
射影次元$22$、一般階数$4$。
源と写像は第\ref{source:circuits51}節の有理演算源$U_i$、$\rho_i$。演算列は\path{reports/ADDITIONAL_RANK4_S61_CIRCUITS_20261103_v62/R4_TRIPLE_HYPERPLANE_QUADRIC.json.gz}。
入力は\path{scale,base1,base2,base3,base4,g0,g1,g2,g3,g4,g5,g6,g7,g8,g9}。全inverse節点と記載した分母の非零開を取る。
一般実退化因子は$H_{1}^{3}+D_{2,2}$。
description：3H+Q; Q has rank 4。
構成は第\ref{source:circuits51}節、成分性は第\ref{sec:known-component-certificates}節、相違は第\ref{sec:known-component-certificates}節及びその証明を参照する。

\paragraph{族49}\label{catalogue:R4_CA}
\path{R4_CA}。
射影次元$28$、一般階数$4$。
源と写像は第\ref{source:circuits51}節の有理演算源$U_i$、$\rho_i$。演算列は\path{reports/ADDITIONAL_RANK4_S61_CIRCUITS_20261103_v62/R4_CA.json.gz}。
入力は\path{a0,a1,a2,b0,b1,b2,b3,b4,b5,b6,b7,b8,b9,b10,b11,b12,b13,b14,b15,scale,base1,base2,base3,base4,g0,g1,g2,g3,g4,g5,g6,g7,g8,g9}。全inverse節点と記載した分母の非零開を取る。
一般実退化因子は$H_{1}^{3}+H_{2}+H_{3}$。
description：3H+H1+H2; three distinct hyperplanes in one pencil。
構成は第\ref{source:circuits51}節、成分性は第\ref{sec:known-component-certificates}節、相違は第\ref{sec:known-component-certificates}節及びその証明を参照する。

\paragraph{族50}\label{catalogue:R4_CB}
\path{R4_CB}。
射影次元$28$、一般階数$4$。
源と写像は第\ref{source:circuits51}節の有理演算源$U_i$、$\rho_i$。演算列は\path{reports/ADDITIONAL_RANK4_S61_CIRCUITS_20261103_v62/R4_CB.json.gz}。
入力は\path{a0,a1,a2,b0,b1,b2,b3,b4,b5,b6,b7,b8,b9,b10,b11,b12,b13,b14,b15,scale,base1,base2,base3,base4,g0,g1,g2,g3,g4,g5,g6,g7,g8,g9}。全inverse節点と記載した分母の非零開を取る。
一般実退化因子は$H_{1}+H_{2}^{2}+H_{3}^{2}$。
description：H+2H1+2H2; three distinct hyperplanes in one pencil。
構成は第\ref{source:circuits51}節、成分性は第\ref{sec:known-component-certificates}節、相違は第\ref{sec:known-component-certificates}節及びその証明を参照する。

\paragraph{族51}\label{catalogue:R4_H_QUARTIC_ESSENTIAL5}
\path{R4_H_QUARTIC_ESSENTIAL5}。
射影次元$23$、一般階数$4$。
源と写像は第\ref{source:circuits51}節の有理演算源$U_i$、$\rho_i$。演算列は\path{reports/ADDITIONAL_RANK4_S61_CIRCUITS_20261103_v62/R4_H_QUARTIC_ESSENTIAL5.json.gz}。
入力は\path{scale,base1,base2,base3,base4,g0,g1,g2,g3,g4,g5,g6,g7,g8,g9}。全inverse節点と記載した分母の非零開を取る。
一般実退化因子は$H_{1}+D_{4,2}$。
description：H+F4; irreducible quartic with five essential variables。
構成は第\ref{source:circuits51}節、成分性は第\ref{sec:known-component-certificates}節、相違は第\ref{sec:known-component-certificates}節及びその証明を参照する。

\paragraph{族52}\label{catalogue:R4_ZERO_LOG_5H}
\path{R4_ZERO_LOG_5H}。
射影次元$23$、一般階数$4$。
源と写像は第\ref{source:circuits51}節の有理演算源$U_i$、$\rho_i$。演算列は\path{reports/ADDITIONAL_RANK4_S61_CIRCUITS_20261103_v62/R4_ZERO_LOG_5H.json.gz}。
入力は\path{a,b,scale,base1,base2,base3,base4,g0,g1,g2,g3,g4,g5,g6,g7,g8,g9}。全inverse節点と記載した分母の非零開を取る。
一般実退化因子は$H_{1}^{5}$。
description：5H。
構成は第\ref{source:circuits51}節、成分性は第\ref{sec:known-component-certificates}節、相違は第\ref{sec:known-component-certificates}節及びその証明を参照する。

\paragraph{族53}\label{catalogue:R4_MINIMAL_M}
\path{R4_MINIMAL_M}。
射影次元$23$、一般階数$4$。
源と写像は第\ref{source:circuits51}節の有理演算源$U_i$、$\rho_i$。演算列は\path{reports/ADDITIONAL_RANK4_S61_CIRCUITS_20261103_v62/R4_MINIMAL_M.json.gz}。
入力は\path{scale,base1,base2,base3,base4,g0,g1,g2,g3,g4,g5,g6,g7,g8,g9}。全inverse節点と記載した分母の非零開を取る。
一般実退化因子は$H_{1}+D_{4,2}$。
description：H+F4; irreducible quartic cone with one singular plane, transverse cubic triple line。
構成は第\ref{source:circuits51}節、成分性は第\ref{sec:known-component-certificates}節、相違は第\ref{sec:known-component-certificates}節及びその証明を参照する。

\paragraph{族54}\label{catalogue:R4_A0}
\path{R4_A0}。
射影次元$23$、一般階数$4$。
源と写像は第\ref{source:circuits51}節の有理演算源$U_i$、$\rho_i$。演算列は\path{reports/ADDITIONAL_RANK4_S61_CIRCUITS_20261103_v62/R4_A0.json.gz}。
入力は\path{scale,base1,base2,base3,base4,g0,g1,g2,g3,g4,g5,g6,g7,g8,g9}。全inverse節点と記載した分母の非零開を取る。
一般実退化因子は$H_{1}+D_{4,2}$。
description：H+F4; irreducible quartic cone with one singular plane, transverse cubic type (2,1)。
構成は第\ref{source:circuits51}節、成分性は第\ref{sec:known-component-certificates}節、相違は第\ref{sec:known-component-certificates}節及びその証明を参照する。

\paragraph{族55}\label{catalogue:R4_DPM}
\path{R4_DPM}。
射影次元$23$、一般階数$4$。
源と写像は第\ref{source:circuits51}節の有理演算源$U_i$、$\rho_i$。演算列は\path{reports/ADDITIONAL_RANK4_S61_CIRCUITS_20261103_v62/R4_DPM.json.gz}。
入力は\path{r,scale,base1,base2,base3,base4,g0,g1,g2,g3,g4,g5,g6,g7,g8,g9}。全inverse節点と記載した分母の非零開を取る。
一般実退化因子は$H_{1}+D_{4,2}$。
description：H+F4; irreducible quartic cone with three singular planes, generic A2 transverse cusps。
構成は第\ref{source:circuits51}節、成分性は第\ref{sec:known-component-certificates}節、相違は第\ref{sec:known-component-certificates}節及びその証明を参照する。

\paragraph{族56}\label{catalogue:R4_CPLUS_EXT}
\path{R4_CPLUS_EXT}。
射影次元$24$、一般階数$4$。
源と写像は第\ref{source:circuits51}節の有理演算源$U_i$、$\rho_i$。演算列は\path{reports/ADDITIONAL_RANK4_S61_CIRCUITS_20261103_v62/R4_CPLUS_EXT.json.gz}。
入力は\path{c,beta,scale,base1,base2,base3,base4,g0,g1,g2,g3,g4,g5,g6,g7,g8,g9}。全inverse節点と記載した分母の非零開を取る。
一般実退化因子は$H_{1}+D_{2,2}+D_{2,3}$。
description：H+Q1+Q2; distinct irreducible quadrics and at most three essential variables in the quintic。
構成は第\ref{source:circuits51}節、成分性は第\ref{sec:known-component-certificates}節、相違は第\ref{sec:known-component-certificates}節及びその証明を参照する。

\paragraph{族57}\label{catalogue:R4_CMINUS_MAG}
\path{R4_CMINUS_MAG}。
射影次元$23$、一般階数$4$。
源と写像は第\ref{source:circuits51}節の有理演算源$U_i$、$\rho_i$。演算列は\path{reports/ADDITIONAL_RANK4_S61_CIRCUITS_20261103_v62/R4_CMINUS_MAG.json.gz}。
入力は\path{beta,gamma,scale,base1,base2,base3,base4,g0,g1,g2,g3,g4,g5,g6,g7,g8,g9}。全inverse節点と記載した分母の非零開を取る。
一般実退化因子は$H_{1}+H_{2}+H_{3}+D_{2,4}$。
description：t*(t\textasciicircum{}2+x\textasciicircum{}2)*(t\textasciicircum{}2+x\textasciicircum{}2-2*z\textasciicircum{}2); three distinct hyperplanes and one rank-three quadric over C。
構成は第\ref{source:circuits51}節、成分性は第\ref{sec:known-component-certificates}節、相違は第\ref{sec:known-component-certificates}節及びその証明を参照する。

\paragraph{族58}\label{catalogue:R4_POSITIVE_WEIGHT6_25}
\path{R4_POSITIVE_WEIGHT6_25}。
射影次元$25$、一般階数$4$。
一般族は命題又は構成\ref{source:new25}の全sourceを用いる。
任意$(n_0,n_1,n_2)\in\mathbb A^3$、一般rank四の開。
一般実退化因子は$H_{1}^{3}+D_{2,2}$。
description：3H+Q, Q smooth rank5 and H tangent to Q。
構成は第\ref{source:new25}節、成分性は第\ref{sec:new25-component}節、相違は第\ref{sec:new25-component}節及びその証明を参照する。

\paragraph{族59}\label{catalogue:R4_DOUBLE_RANK3_QUADRIC_MIXED_22}
\path{R4_DOUBLE_RANK3_QUADRIC_MIXED_22}。
射影次元$22$、一般階数$4$。
一般族は命題又は構成\ref{prop:new22}の全sourceを用いる。
一般実退化因子は$H_{1}+D_{2,2}^{2}$。
quadric\_rank：3。
description：H+2Q; Q has rank 3 and H contains the vertex and is tangent to the quotient conic。
構成は第\ref{prop:new22}節、成分性は第\ref{prop:new22}節、相違は第\ref{prop:new22}節及びその証明を参照する。

\paragraph{族60}\label{catalogue:R4_DOUBLE_RANK3_QUADRIC_TWZERO_24}
\path{R4_DOUBLE_RANK3_QUADRIC_TWZERO_24}。
射影次元$24$、一般階数$4$。
一般族は命題又は構成\ref{prop:new24}の全sourceを用いる。
\path{source}："SL2 x Gm x A7"。
一般実退化因子は$H_{1}+D_{2,2}^{2}$。
quadric\_rank：3。
description：H+2Q; Q has rank3 and H contains the vertex and is tangent to the quotient conic。
構成は第\ref{prop:new24}節、成分性は第\ref{prop:new24}節、相違は第\ref{prop:new24}節及びその証明を参照する。

\paragraph{族61}\label{catalogue:R4_NIL4_DOUBLE_H_CUBIC_22}
\path{R4_NIL4_DOUBLE_H_CUBIC_22}。
射影次元$22$、一般階数$4$。
一般族は命題又は構成\ref{prop:nil-four-components}の全sourceを用いる。
枝 \path{source_parameter_s="1/3"}。
一般実退化因子は$H_{1}^{2}+D_{3,2}$。
description：冪零Jordan4・二重超平面と既約三次型。
構成は第\ref{prop:nil-four-components}節、成分性は第\ref{prop:nil-four-components}節、相違は第\ref{prop:nil-four-components}節及びその証明を参照する。

\paragraph{族62}\label{catalogue:R4_NIL4_DOUBLE_H_H_QUADRIC_22}
\path{R4_NIL4_DOUBLE_H_H_QUADRIC_22}。
射影次元$22$、一般階数$4$。
一般族は命題又は構成\ref{prop:nil-four-components}の全sourceを用いる。
枝 \path{source_parameter_s="1/7"}。
一般実退化因子は$H_{1}^{2}+H_{2}+D_{2,3}$。
quadric\_rank：3。
description：冪零Jordan4・二重超平面と超平面と階数三二次型。
構成は第\ref{prop:nil-four-components}節、成分性は第\ref{prop:nil-four-components}節、相違は第\ref{prop:nil-four-components}節及びその証明を参照する。

\paragraph{族63}\label{catalogue:R4_NIL4_TRIPLE_H_DOUBLE_H_OUTSIDE_22}
\path{R4_NIL4_TRIPLE_H_DOUBLE_H_OUTSIDE_22}。
射影次元$22$、一般階数$4$。
一般族は命題又は構成\ref{prop:nil-four-components}の全sourceを用いる。
枝 \path{source_parameter_s="1/5"}。
一般実退化因子は$H_{1}^{3}+H_{2}^{2}$。
description：三重超平面cofactor像の外にあるJordan4型。
構成は第\ref{prop:nil-four-components}節、成分性は第\ref{prop:nil-four-components}節、相違は第\ref{prop:nil-four-components}節及びその証明を参照する。

\paragraph{族64}\label{catalogue:R4_NIL4_TRIPLE_H_DOUBLE_H_INSIDE_22}
\path{R4_NIL4_TRIPLE_H_DOUBLE_H_INSIDE_22}。
射影次元$22$、一般階数$4$。
一般族は命題又は構成\ref{prop:nil-four-components}の全sourceを用いる。
枝 \path{source_parameter_s="1/15"}。
一般実退化因子は$H_{1}^{3}+H_{2}^{2}$。
description：三重超平面cofactor像に入るJordan4型。
構成は第\ref{prop:nil-four-components}節、成分性は第\ref{prop:nil-four-components}節、相違は第\ref{prop:nil-four-components}節及びその証明を参照する。

\paragraph{族65}\label{catalogue:R4_DEPENDENT_THREE_HYPERPLANES_QUADRIC_24}
\path{R4_DEPENDENT_THREE_HYPERPLANES_QUADRIC_24}。
射影次元$24$、一般階数$4$。
一般族は命題又は構成\ref{prop:dependent-three-component}の全sourceを用いる。
\path{source_open}："a0*(a0+a1) != 0; R arbitrary polynomial of degree at most2"。
一般実退化因子は$H_{1}+H_{2}+H_{3}+D_{2,4}$。
quadric\_rank：3。
description：three distinct dependent hyperplanes and one irreducible rank3 quadric。
構成は第\ref{prop:dependent-three-component}節、成分性は第\ref{prop:dependent-three-component}節、相違は第\ref{prop:dependent-three-component}節及びその証明を参照する。

\paragraph{族66}\label{catalogue:R4_MODULAR_NIL31_DOUBLE_DOUBLE_H_25}
\path{R4_MODULAR_NIL31_DOUBLE_DOUBLE_H_25}。
射影次元$25$、一般階数$4$。
一般族は命題又は構成\ref{prop:nil31-component}の全sourceを用いる。
\path{source}："A9 x Gm"。
一般実退化因子は$H_{1}^{2}+H_{2}^{2}+H_{3}$。
description：three independent hyperplanes with multiplicities2,2,1。
構成は第\ref{prop:nil31-component}節、成分性は第\ref{prop:nil31-component}節、相違は第\ref{prop:nil31-component}節及びその証明を参照する。

\paragraph{族67}\label{catalogue:R4_NIL31_SHORT_DOUBLE_H_CUBIC_22}
\path{R4_NIL31_SHORT_DOUBLE_H_CUBIC_22}。
射影次元$22$、一般階数$4$。
一般族は命題又は構成\ref{prop:nil31-three22}の全sourceを用いる。
枝 \path{source_branch="r=-1/3"}。
一般実退化因子は$H_{1}^{2}+D_{3,2}$。
description：冪零短鎖定数型・二重超平面と既約三次成分。
構成は第\ref{prop:nil31-three22}節、成分性は第\ref{prop:nil31-three22}節、相違は第\ref{prop:nil31-three22}節及びその証明を参照する。

\paragraph{族68}\label{catalogue:R4_NIL31_LONG_DOUBLE_H_CUBIC_22}
\path{R4_NIL31_LONG_DOUBLE_H_CUBIC_22}。
射影次元$22$、一般階数$4$。
一般族は命題又は構成\ref{prop:nil31-three22}の全sourceを用いる。
枝 \path{source_branch="E: r=2/3"}。
一般実退化因子は$H_{1}^{2}+D_{3,2}$。
description：冪零長鎖定数型・二重超平面と既約三次成分。
構成は第\ref{prop:nil31-three22}節、成分性は第\ref{prop:nil31-three22}節、相違は第\ref{prop:nil31-three22}節及びその証明を参照する。

\paragraph{族69}\label{catalogue:R4_NIL31_LONG_TRIPLE_H_RANK3_QUADRIC_22}
\path{R4_NIL31_LONG_TRIPLE_H_RANK3_QUADRIC_22}。
射影次元$22$、一般階数$4$。
一般族は命題又は構成\ref{prop:nil31-three22}の全sourceを用いる。
枝 \path{source_branch="F: r=3/2"}。
一般実退化因子は$H_{1}^{3}+D_{2,2}$。
quadric\_rank：3。
description：冪零長鎖定数型・三重超平面と階数三二次成分。
構成は第\ref{prop:nil31-three22}節、成分性は第\ref{prop:nil31-three22}節、相違は第\ref{prop:nil31-three22}節及びその証明を参照する。

\paragraph{族70}\label{catalogue:R4_INDEPENDENT_THREE_HYPERPLANES_TRITANGENT_QUADRIC_23}
\path{R4_INDEPENDENT_THREE_HYPERPLANES_TRITANGENT_QUADRIC_23}。
射影次元$23$、一般階数$4$。
一般族は命題又は構成\ref{prop:independent-three-component}の全sourceを用いる。
一般実退化因子は$H_{1}+H_{2}+H_{3}+D_{2,4}$。
quadric\_rank：3。
essential\_variable\_count：3。
description：three independent hyperplanes and a rank3 quadric tangent to all three coordinate sides。
構成は第\ref{prop:independent-three-component}節、成分性は第\ref{prop:independent-three-component}節、相違は第\ref{prop:independent-three-component}節及びその証明を参照する。

\paragraph{族71}\label{catalogue:R4_NIL22_QUADRUPLE_H_LINEAR_22}
\path{R4_NIL22_QUADRUPLE_H_LINEAR_22}。
射影次元$22$、一般階数$4$。
一般族は命題又は構成\ref{prop:weighted-three22}の全sourceを用いる。
枝 \path{source_branch="A=-1"}。
一般実退化因子は$H_{1}^{4}+H_{2}$。
description：nil22定数型・四重超平面と一次因子。
構成は第\ref{prop:weighted-three22}節、成分性は第\ref{prop:weighted-three22}節、相違は第\ref{prop:weighted-three22}節及びその証明を参照する。

\paragraph{族72}\label{catalogue:R4_NIL31_SHORT_QUADRUPLE_H_LINEAR_22}
\path{R4_NIL31_SHORT_QUADRUPLE_H_LINEAR_22}。
射影次元$22$、一般階数$4$。
一般族は命題又は構成\ref{prop:weighted-three22}の全sourceを用いる。
枝 \path{source_branch="r=1"}。
一般実退化因子は$H_{1}^{4}+H_{2}$。
description：nil31短鎖型・四重超平面と一次因子。
構成は第\ref{prop:weighted-three22}節、成分性は第\ref{prop:weighted-three22}節、相違は第\ref{prop:weighted-three22}節及びその証明を参照する。

\paragraph{族73}\label{catalogue:R4_NIL22_DOUBLE_H_TRIPLE_LINEAR_22}
\path{R4_NIL22_DOUBLE_H_TRIPLE_LINEAR_22}。
射影次元$22$、一般階数$4$。
一般族は命題又は構成\ref{prop:weighted-three22}の全sourceを用いる。
枝 \path{source_branch="A=-1/3, equivalent nil31 short r=1/2"}。
一般実退化因子は$H_{1}^{2}+H_{2}^{3}$。
description：nil22共鳴型・二重超平面と三重一次因子。
構成は第\ref{prop:weighted-three22}節、成分性は第\ref{prop:weighted-three22}節、相違は第\ref{prop:weighted-three22}節及びその証明を参照する。

\paragraph{族74}\label{catalogue:R4_TWO_HYPERPLANES_CUBIC_22}
\path{R4_TWO_HYPERPLANES_CUBIC_22}。
射影次元$22$、一般階数$4$。
一般族は命題又は構成\ref{prop:two-h-cubic22}の全sourceを用いる。
枝 \path{source_branches=["p=1", "p=-1"]}。
一般実退化因子は$H_{1}+H_{2}+D_{3,3}$。
description：two distinct hyperplanes and an irreducible cubic。
構成は第\ref{prop:two-h-cubic22}節、成分性は第\ref{prop:two-h-cubic22}節、相違は第\ref{prop:two-h-cubic22}節及びその証明を参照する。

\paragraph{族75}\label{catalogue:R4_PLANE_CONIC_CUBIC_CUSP_OUTSIDE_1_23}
\path{R4_PLANE_CONIC_CUBIC_CUSP_OUTSIDE_1_23}。
射影次元$23$、一般階数$4$。
一般族は命題又は構成\ref{prop:plane-conic-cubic-six}の全sourceを用いる。
枝 \path{source_case="cusp_outside_1"}。
磁気項は任意二次式$R_2$の全六係数である。
一般実退化因子は$D_{2,1}+D_{3,2}$。
quadric\_rank：3。
essential\_variable\_count：3。
構成は第\ref{prop:plane-conic-cubic-six}節、成分性は第\ref{prop:plane-conic-cubic-six}節、相違は第\ref{prop:plane-conic-cubic-six}節及びその証明を参照する。

\paragraph{族76}\label{catalogue:R4_PLANE_CONIC_CUBIC_CUSP_OUTSIDE_3_23}
\path{R4_PLANE_CONIC_CUBIC_CUSP_OUTSIDE_3_23}。
射影次元$23$、一般階数$4$。
一般族は命題又は構成\ref{prop:plane-conic-cubic-six}の全sourceを用いる。
枝 \path{source_case="cusp_outside_3"}。
磁気項は任意二次式$R_2$の全六係数である。
一般実退化因子は$D_{2,1}+D_{3,2}$。
quadric\_rank：3。
essential\_variable\_count：3。
構成は第\ref{prop:plane-conic-cubic-six}節、成分性は第\ref{prop:plane-conic-cubic-six}節、相違は第\ref{prop:plane-conic-cubic-six}節及びその証明を参照する。

\paragraph{族77}\label{catalogue:R4_PLANE_CONIC_CUBIC_CUSP_AT_CUSP_2_23}
\path{R4_PLANE_CONIC_CUBIC_CUSP_AT_CUSP_2_23}。
射影次元$23$、一般階数$4$。
一般族は命題又は構成\ref{prop:plane-conic-cubic-six}の全sourceを用いる。
枝 \path{source_case="cusp_at_cusp_2"}。
磁気項は任意二次式$R_2$の全六係数である。
一般実退化因子は$D_{2,1}+D_{3,2}$。
quadric\_rank：3。
essential\_variable\_count：3。
構成は第\ref{prop:plane-conic-cubic-six}節、成分性は第\ref{prop:plane-conic-cubic-six}節、相違は第\ref{prop:plane-conic-cubic-six}節及びその証明を参照する。

\paragraph{族78}\label{catalogue:R4_PLANE_CONIC_CUBIC_CUSP_AT_CUSP_3_23}
\path{R4_PLANE_CONIC_CUBIC_CUSP_AT_CUSP_3_23}。
射影次元$23$、一般階数$4$。
一般族は命題又は構成\ref{prop:plane-conic-cubic-six}の全sourceを用いる。
枝 \path{source_case="cusp_at_cusp_3"}。
磁気項は任意二次式$R_2$の全六係数である。
一般実退化因子は$D_{2,1}+D_{3,2}$。
quadric\_rank：3。
essential\_variable\_count：3。
構成は第\ref{prop:plane-conic-cubic-six}節、成分性は第\ref{prop:plane-conic-cubic-six}節、相違は第\ref{prop:plane-conic-cubic-six}節及びその証明を参照する。

\paragraph{族79}\label{catalogue:R4_PLANE_CONIC_CUBIC_NODE_SMOOTH_CONTACT_23}
\path{R4_PLANE_CONIC_CUBIC_NODE_SMOOTH_CONTACT_23}。
射影次元$23$、一般階数$4$。
一般族は命題又は構成\ref{prop:plane-conic-cubic-six}の全sourceを用いる。
枝 \path{source_case="node_smooth_contact"}。
磁気項は任意二次式$R_2$の全六係数である。
一般実退化因子は$D_{2,1}+D_{3,2}$。
quadric\_rank：3。
essential\_variable\_count：3。
構成は第\ref{prop:plane-conic-cubic-six}節、成分性は第\ref{prop:plane-conic-cubic-six}節、相違は第\ref{prop:plane-conic-cubic-six}節及びその証明を参照する。

\paragraph{族80}\label{catalogue:R4_PLANE_CONIC_CUBIC_NODE_AT_NODE_23}
\path{R4_PLANE_CONIC_CUBIC_NODE_AT_NODE_23}。
射影次元$23$、一般階数$4$。
一般族は命題又は構成\ref{prop:plane-conic-cubic-six}の全sourceを用いる。
枝 \path{source_case="node_at_node"}。
磁気項は任意二次式$R_2$の全六係数である。
一般実退化因子は$D_{2,1}+D_{3,2}$。
quadric\_rank：3。
essential\_variable\_count：3。
構成は第\ref{prop:plane-conic-cubic-six}節、成分性は第\ref{prop:plane-conic-cubic-six}節、相違は第\ref{prop:plane-conic-cubic-six}節及びその証明を参照する。

\paragraph{族81}\label{catalogue:R4_PLANE_TWO_LINES_CUBIC_NODE_TANGENT_FLEX_23}
\path{R4_PLANE_TWO_LINES_CUBIC_NODE_TANGENT_FLEX_23}。
射影次元$23$、一般階数$4$。
一般族は命題又は構成\ref{prop:plane-two-lines-cubic-six}の全sourceを用いる。
枝 \path{source_case="node_tangent_flex"}。
磁気項は任意二次式$R_2$の全六係数である。
一般実退化因子は$H_{1}+H_{2}+D_{3,3}$。
essential\_variable\_count：3。
構成は第\ref{prop:plane-two-lines-cubic-six}節、成分性は第\ref{prop:plane-two-lines-cubic-six}節、相違は第\ref{prop:plane-two-lines-cubic-six}節及びその証明を参照する。

\paragraph{族82}\label{catalogue:R4_PLANE_TWO_LINES_CUBIC_CUSP_FLEX_TANGENT_23}
\path{R4_PLANE_TWO_LINES_CUBIC_CUSP_FLEX_TANGENT_23}。
射影次元$23$、一般階数$4$。
一般族は命題又は構成\ref{prop:plane-two-lines-cubic-six}の全sourceを用いる。
枝 \path{source_case="cusp_flex_tangent"}。
磁気項は任意二次式$R_2$の全六係数である。
一般実退化因子は$H_{1}+H_{2}+D_{3,3}$。
essential\_variable\_count：3。
構成は第\ref{prop:plane-two-lines-cubic-six}節、成分性は第\ref{prop:plane-two-lines-cubic-six}節、相違は第\ref{prop:plane-two-lines-cubic-six}節及びその証明を参照する。

\paragraph{族83}\label{catalogue:R4_PLANE_TWO_LINES_CUBIC_CUSP_TWO_TANGENTS_INCIDENT_23}
\path{R4_PLANE_TWO_LINES_CUBIC_CUSP_TWO_TANGENTS_INCIDENT_23}。
射影次元$23$、一般階数$4$。
一般族は命題又は構成\ref{prop:plane-two-lines-cubic-six}の全sourceを用いる。
枝 \path{source_case="cusp_two_tangents_incident"}。
磁気項は任意二次式$R_2$の全六係数である。
一般実退化因子は$H_{1}+H_{2}+D_{3,3}$。
essential\_variable\_count：3。
構成は第\ref{prop:plane-two-lines-cubic-six}節、成分性は第\ref{prop:plane-two-lines-cubic-six}節、相違は第\ref{prop:plane-two-lines-cubic-six}節及びその証明を参照する。

\paragraph{族84}\label{catalogue:R4_PLANE_TWO_LINES_CUBIC_CUSP_FLEX_CUSP_SECANT_23}
\path{R4_PLANE_TWO_LINES_CUBIC_CUSP_FLEX_CUSP_SECANT_23}。
射影次元$23$、一般階数$4$。
一般族は命題又は構成\ref{prop:plane-two-lines-cubic-six}の全sourceを用いる。
枝 \path{source_case="cusp_flex_cusp_secant"}。
磁気項は任意二次式$R_2$の全六係数である。
一般実退化因子は$H_{1}+H_{2}+D_{3,3}$。
essential\_variable\_count：3。
構成は第\ref{prop:plane-two-lines-cubic-six}節、成分性は第\ref{prop:plane-two-lines-cubic-six}節、相違は第\ref{prop:plane-two-lines-cubic-six}節及びその証明を参照する。

\paragraph{族85}\label{catalogue:R4_PLANE_TWO_LINES_CUBIC_CUSP_TANGENT_CUSP_SECANT_23}
\path{R4_PLANE_TWO_LINES_CUBIC_CUSP_TANGENT_CUSP_SECANT_23}。
射影次元$23$、一般階数$4$。
一般族は命題又は構成\ref{prop:plane-two-lines-cubic-six}の全sourceを用いる。
枝 \path{source_case="cusp_tangent_cusp_secant"}。
磁気項は任意二次式$R_2$の全六係数である。
一般実退化因子は$H_{1}+H_{2}+D_{3,3}$。
essential\_variable\_count：3。
構成は第\ref{prop:plane-two-lines-cubic-six}節、成分性は第\ref{prop:plane-two-lines-cubic-six}節、相違は第\ref{prop:plane-two-lines-cubic-six}節及びその証明を参照する。

\paragraph{族86}\label{catalogue:R4_PLANE_TWO_LINES_CUBIC_CUSP_CUSP_TANGENT_ORDINARY_TANGENT_23}
\path{R4_PLANE_TWO_LINES_CUBIC_CUSP_CUSP_TANGENT_ORDINARY_TANGENT_23}。
射影次元$23$、一般階数$4$。
一般族は命題又は構成\ref{prop:plane-two-lines-cubic-six}の全sourceを用いる。
枝 \path{source_case="cusp_cusp_tangent_ordinary_tangent"}。
磁気項は任意二次式$R_2$の全六係数である。
一般実退化因子は$H_{1}+H_{2}+D_{3,3}$。
essential\_variable\_count：3。
構成は第\ref{prop:plane-two-lines-cubic-six}節、成分性は第\ref{prop:plane-two-lines-cubic-six}節、相違は第\ref{prop:plane-two-lines-cubic-six}節及びその証明を参照する。

\paragraph{族87}\label{catalogue:R4_PLANE_TWO_LINES_CUBIC_NODE_TWO_FLEXES_23}
\path{R4_PLANE_TWO_LINES_CUBIC_NODE_TWO_FLEXES_23}。
射影次元$23$、一般階数$4$。
一般族は命題又は構成\ref{prop:plane-two-lines-cubic-three}の全sourceを用いる。
枝 \path{source_case="node_two_flexes"}。
磁気項は任意二次式$R_2$の全六係数である。
一般実退化因子は$H_{1}+H_{2}+D_{3,3}$。
essential\_variable\_count：3。
構成は第\ref{prop:plane-two-lines-cubic-three}節、成分性は第\ref{prop:plane-two-lines-cubic-three}節、相違は第\ref{prop:plane-two-lines-cubic-three}節及びその証明を参照する。

\paragraph{族88}\label{catalogue:R4_PLANE_TWO_LINES_CUBIC_NODE_FLEX_SECANT_23}
\path{R4_PLANE_TWO_LINES_CUBIC_NODE_FLEX_SECANT_23}。
射影次元$23$、一般階数$4$。
一般族は命題又は構成\ref{prop:plane-two-lines-cubic-three}の全sourceを用いる。
枝 \path{source_case="node_flex_secant"}。
磁気項は任意二次式$R_2$の全六係数である。
一般実退化因子は$H_{1}+H_{2}+D_{3,3}$。
essential\_variable\_count：3。
構成は第\ref{prop:plane-two-lines-cubic-three}節、成分性は第\ref{prop:plane-two-lines-cubic-three}節、相違は第\ref{prop:plane-two-lines-cubic-three}節及びその証明を参照する。

\paragraph{族89}\label{catalogue:R4_PLANE_TWO_LINES_CUBIC_CUSP_TANGENT_CUSP_THIRD_SECANT_23}
\path{R4_PLANE_TWO_LINES_CUBIC_CUSP_TANGENT_CUSP_THIRD_SECANT_23}。
射影次元$23$、一般階数$4$。
一般族は命題又は構成\ref{prop:plane-two-lines-cubic-three}の全sourceを用いる。
枝 \path{source_case="cusp_tangent_cusp_third_secant"}。
磁気項は任意二次式$R_2$の全六係数である。
一般実退化因子は$H_{1}+H_{2}+D_{3,3}$。
essential\_variable\_count：3。
構成は第\ref{prop:plane-two-lines-cubic-three}節、成分性は第\ref{prop:plane-two-lines-cubic-three}節、相違は第\ref{prop:plane-two-lines-cubic-three}節及びその証明を参照する。

\paragraph{族90}\label{catalogue:R4_NIL4_HYPERPLANE_QUARTIC_22}
\path{R4_NIL4_HYPERPLANE_QUARTIC_22}。
射影次元$22$、一般階数$4$。
一般族は命題又は構成\ref{prop:nil4-quartic-whole-first-source}の全sourceを用いる。
枝 \path{source_branch="nil4,c=0, normalized delta=-1/3,n1!=0"}。
一般実退化因子は$H_{1}+D_{4,2}$。
description：simple hyperplane and irreducible quartic。
構成は第\ref{prop:nil4-quartic-whole-first-source}節、成分性は第\ref{prop:nil4-hyperplane-quartic-component}節、相違は第\ref{prop:nil4-quartic-distinctness}節及びその証明を参照する。

\paragraph{族91}\label{catalogue:R4_PLANE_TWO_CONICS_BITANGENT_23}
\path{R4_PLANE_TWO_CONICS_BITANGENT_23}。
射影次元$23$、一般階数$4$。
一般族は命題又は構成\ref{prop:plane-two-conics-two-components}の全sourceを用いる。
枝 \path{configuration="contact31_common_bitangent"}。
磁気項は任意二次式$R_2$の全六係数である。
一般実退化因子は$H_{1}+D_{2,2}+D_{2,3}$。
description：line and two distinct irreducible plane conics, pulled back to P4。
構成は第\ref{prop:plane-two-conics-two-components}節、成分性は第\ref{prop:plane-two-conics-two-components}節、相違は第\ref{prop:plane-two-conics-two-components}節及びその証明を参照する。

\paragraph{族92}\label{catalogue:R4_PLANE_TWO_CONICS_CONTACT4_TANGENT_23}
\path{R4_PLANE_TWO_CONICS_CONTACT4_TANGENT_23}。
射影次元$23$、一般階数$4$。
一般族は命題又は構成\ref{prop:plane-two-conics-two-components}の全sourceを用いる。
枝 \path{configuration="contact4_other_tangent"}。
磁気項は任意二次式$R_2$の全六係数である。
一般実退化因子は$H_{1}+D_{2,2}+D_{2,3}$。
description：line and two distinct irreducible plane conics, pulled back to P4。
構成は第\ref{prop:plane-two-conics-two-components}節、成分性は第\ref{prop:plane-two-conics-two-components}節、相違は第\ref{prop:plane-two-conics-two-components}節及びその証明を参照する。

\paragraph{族93}\label{catalogue:R4_NIL31_IMAGE_HALF_25}
\path{R4_NIL31_IMAGE_HALF_25}。
射影次元$25$、一般階数$4$。
一般族は命題又は構成\ref{prop:nil31-image-two-components}の全sourceを用いる。
枝 \path{branch="half"}。
枝 \path{fixed_r="-1/2"}。
\path{free_parameters}：["A0", "A1", "A2", "A3", "A4", "A5", "u", "v", "t", "p"]。
一般実退化因子は$H_{1}^{2}+H_{2}^{3}$。
description：Two distinct hyperplanes with multiplicities three and two。
構成は第\ref{prop:nil31-image-two-components}節、成分性は第\ref{prop:nil31-image-two-components}節、相違は第\ref{prop:nil31-image-two-components}節及びその証明を参照する。

\paragraph{族94}\label{catalogue:R4_NIL31_IMAGE_THIRD_25}
\path{R4_NIL31_IMAGE_THIRD_25}。
射影次元$25$、一般階数$4$。
一般族は命題又は構成\ref{prop:nil31-image-two-components}の全sourceを用いる。
枝 \path{branch="third"}。
枝 \path{fixed_r="-1/3"}。
\path{free_parameters}：["A0", "A1", "A2", "A3", "A4", "A5", "u", "v", "t", "q"]。
一般実退化因子は$H_{1}^{3}+H_{2}^{2}$。
description：Two distinct hyperplanes with multiplicities three and two。
構成は第\ref{prop:nil31-image-two-components}節、成分性は第\ref{prop:nil31-image-two-components}節、相違は第\ref{prop:nil31-image-two-components}節及びその証明を参照する。

\paragraph{族95}\label{catalogue:R4_NIL22_IMAGE_NODAL_CUBIC_FLEX_23}
\path{R4_NIL22_IMAGE_NODAL_CUBIC_FLEX_23}。
射影次元$23$、一般階数$4$。
一般族は命題又は構成\ref{prop:nil22-image-two-components}の全sourceを用いる。
枝 \path{branch="triple_A"}。
磁気項は任意二次式$R_2$の全六係数である。
一般実退化因子は$H_{1}^{2}+D_{3,2}$。
description：Unique double hyperplane plus irreducible nodal cubic cone。
構成は第\ref{prop:nil22-image-two-components}節、成分性は第\ref{prop:nil22-image-two-components}節、相違は第\ref{prop:nil22-image-two-components}節及びその証明を参照する。

\paragraph{族96}\label{catalogue:R4_NIL22_IMAGE_NODAL_CUBIC_NODE_TANGENT_23}
\path{R4_NIL22_IMAGE_NODAL_CUBIC_NODE_TANGENT_23}。
射影次元$23$、一般階数$4$。
一般族は命題又は構成\ref{prop:nil22-image-two-components}の全sourceを用いる。
枝 \path{branch="triple_B"}。
磁気項は任意二次式$R_2$の全六係数である。
一般実退化因子は$H_{1}^{2}+D_{3,2}$。
description：Unique double hyperplane plus irreducible nodal cubic cone。
構成は第\ref{prop:nil22-image-two-components}節、成分性は第\ref{prop:nil22-image-two-components}節、相違は第\ref{prop:nil22-image-two-components}節及びその証明を参照する。

\paragraph{族97}\label{catalogue:R4_NIL31_LONG_HYPERPLANE_QUARTIC_22}
\path{R4_NIL31_LONG_HYPERPLANE_QUARTIC_22}。
射影次元$22$、一般階数$4$。
一般族は命題又は構成\ref{prop:nil31-long-quartic-component}の全sourceを用いる。
\path{generic_condition}："a!=0"。
\path{free_parameters}：["a", "n3", "n4", "n8", "n12"]。
一般実退化因子は$H_{1}+D_{4,2}$。
description：Unique simple hyperplane and irreducible binary-cubic discriminant quartic。
構成は第\ref{prop:nil31-long-quartic-component}節、成分性は第\ref{prop:nil31-long-quartic-component}節、相違は第\ref{prop:nil31-long-quartic-component}節及びその証明を参照する。

\paragraph{族98}\label{catalogue:R4_GA9_CONSTANT_HYPERPLANE_QUARTIC_23}
\path{R4_GA9_CONSTANT_HYPERPLANE_QUARTIC_23}。
射影次元$23$、一般階数$4$。
一般族は命題又は構成\ref{prop:ga9-constant-component}の全sourceを用いる。
\path{generic_condition}："r!=0"。
\path{free_parameters}：["r", "b", "alpha", "p", "q", "A", "D", "G", "H", "k"]。
一般実退化因子は$H_{1}+D_{4,2}$。
description：Unique simple hyperplane and absolutely irreducible quartic; quintic vanishes to order at least three on a plane。
構成は第\ref{prop:ga9-constant-component}節、成分性は第\ref{prop:ga9-constant-component}節、相違は第\ref{prop:ga9-constant-component}節及びその証明を参照する。

\paragraph{族99}\label{catalogue:R4_GA9_DISTINCT_IRREDUCIBLE_QUINTIC_23}
\path{R4_GA9_DISTINCT_IRREDUCIBLE_QUINTIC_23}。
射影次元$23$、一般階数$4$。
一般族は命題又は構成\ref{prop:Ga9-distinct-quintic-component}の全sourceを用いる。
一般実退化因子は$D_{5,1}$。
description：Absolutely irreducible reduced nonnormal quintic with two singular planes, generic multiplicities3 and2。
構成は第\ref{prop:Ga9-distinct-quintic-component}節、成分性は第\ref{prop:Ga9-distinct-quintic-component}節、相違は第\ref{prop:Ga9-distinct-quintic-component}節及びその証明を参照する。

\paragraph{族100}\label{catalogue:R4_NIL22_LONG_HYPERPLANE_QUARTIC_CUBIC_22}
\path{R4_NIL22_LONG_HYPERPLANE_QUARTIC_CUBIC_22}。
射影次元$22$、一般階数$4$。
一般族は命題又は構成\ref{prop:nil22-long-two-components}の全sourceを用いる。
枝 \path{source_index=3}。
一般実退化因子は$H_{1}+D_{4,2}$。
description：Simple hyperplane plus absolutely irreducible quartic。
構成は第\ref{prop:nil22-long-two-components}節、成分性は第\ref{prop:nil22-long-two-components}節、相違は第\ref{prop:nil22-long-two-components}節及びその証明を参照する。

\paragraph{族101}\label{catalogue:R4_NIL22_LONG_HYPERPLANE_QUARTIC_QUADRATIC_22}
\path{R4_NIL22_LONG_HYPERPLANE_QUARTIC_QUADRATIC_22}。
射影次元$22$、一般階数$4$。
一般族は命題又は構成\ref{prop:nil22-long-two-components}の全sourceを用いる。
枝 \path{source_index=4}。
一般実退化因子は$H_{1}+D_{4,2}$。
description：Simple hyperplane plus absolutely irreducible quartic。
構成は第\ref{prop:nil22-long-two-components}節、成分性は第\ref{prop:nil22-long-two-components}節、相違は第\ref{prop:nil22-long-two-components}節及びその証明を参照する。

\paragraph{族102}\label{catalogue:R4_NIL22_LONG_IRREDUCIBLE_QUINTIC_23}
\path{R4_NIL22_LONG_IRREDUCIBLE_QUINTIC_23}。
射影次元$23$、一般階数$4$。
一般族は命題又は構成\ref{prop:nil22-long-five-curves}の全sourceを用いる。
一般実退化因子は$D_{5,1}$。
description：Absolutely irreducible reduced nonnormal quintic with three singular planes。
構成は第\ref{prop:nil22-long-five-curves}節、成分性は第\ref{prop:nil22-long-five-curves}節、相違は第\ref{prop:nil22-long-five-curves}節及びその証明を参照する。

\paragraph{族103}\label{catalogue:R4_GA9_CUSPIDAL_PLANE_QUARTIC_24}
\path{R4_GA9_CUSPIDAL_PLANE_QUARTIC_24}。
射影次元$24$、一般階数$4$。
一般族は命題又は構成\ref{prop:Ga9-cuspidal-quartic-component}の全sourceを用いる。
\path{full_source_parameters}：["a", "r0", "r1", "r2", "r3", "r4", "r5"]。
磁気項は任意二次式$R_2$の全六係数である。
一般実退化因子は$H_{1}+D_{4,2}$。
essential\_variable\_count：3。
description：Simple H and absolutely irreducible rational plane quartic with A3 tacnode and A2 cusp; H is the common tacnode tangent with total contact4。
構成は第\ref{prop:Ga9-cuspidal-quartic-component}節、成分性は第\ref{prop:Ga9-cuspidal-quartic-component}節、相違は第\ref{prop:Ga9-cuspidal-quartic-component}節及びその証明を参照する。

\paragraph{族104}\label{catalogue:R4_PLANE_UNICUSPIDAL_QUINTIC_45_23}
\path{R4_PLANE_UNICUSPIDAL_QUINTIC_45_23}。
射影次元$23$、一般階数$4$。
一般族は命題又は構成\ref{prop:plane-unicuspidal-quintic-45-component}の全sourceを用いる。
磁気項は任意二次式$R_2$の全六係数である。
一般実退化因子は$D_{5,1}$。
essential\_variable\_count：3。
description：Absolutely irreducible essential-three plane quintic with a unique weighted homogeneous (4,5) cusp, totalMilnor=Tjurina12。
構成は第\ref{prop:plane-unicuspidal-quintic-45-component}節、成分性は第\ref{prop:plane-unicuspidal-quintic-45-component}節、相違は第\ref{prop:plane-unicuspidal-quintic-45-component}節及びその証明を参照する。

\paragraph{族105}\label{catalogue:R4_PLANE_QUINTIC_A12_23}
\path{R4_PLANE_QUINTIC_A12_23}。
射影次元$23$、一般階数$4$。
一般族は命題又は構成\ref{prop:essential-three-eight-components}の全sourceを用いる。
磁気項は任意二次式$R_2$の全六係数である。
一般実退化因子は$D_{5,1}$。
essential\_variable\_count：3。
description：Absolutely irreducible essential-three rational cuspidal quintic; configuration A12。
構成は第\ref{prop:essential-three-eight-components}節、成分性は第\ref{prop:essential-three-eight-components}節、相違は第\ref{prop:essential-three-eight-components}節及びその証明を参照する。

\paragraph{族106}\label{catalogue:R4_PLANE_QUINTIC_E8_A4_23}
\path{R4_PLANE_QUINTIC_E8_A4_23}。
射影次元$23$、一般階数$4$。
一般族は命題又は構成\ref{prop:essential-three-eight-components}の全sourceを用いる。
磁気項は任意二次式$R_2$の全六係数である。
一般実退化因子は$D_{5,1}$。
essential\_variable\_count：3。
description：Absolutely irreducible essential-three rational cuspidal quintic; configuration E8+A4。
構成は第\ref{prop:essential-three-eight-components}節、成分性は第\ref{prop:essential-three-eight-components}節、相違は第\ref{prop:essential-three-eight-components}節及びその証明を参照する。

\paragraph{族107}\label{catalogue:R4_PLANE_QUINTIC_E6_A6_23}
\path{R4_PLANE_QUINTIC_E6_A6_23}。
射影次元$23$、一般階数$4$。
一般族は命題又は構成\ref{prop:essential-three-eight-components}の全sourceを用いる。
磁気項は任意二次式$R_2$の全六係数である。
一般実退化因子は$D_{5,1}$。
essential\_variable\_count：3。
description：Absolutely irreducible essential-three rational cuspidal quintic; configuration E6+A6。
構成は第\ref{prop:essential-three-eight-components}節、成分性は第\ref{prop:essential-three-eight-components}節、相違は第\ref{prop:essential-three-eight-components}節及びその証明を参照する。

\paragraph{族108}\label{catalogue:R4_PLANE_QUINTIC_E6_A4_A2_23}
\path{R4_PLANE_QUINTIC_E6_A4_A2_23}。
射影次元$23$、一般階数$4$。
一般族は命題又は構成\ref{prop:essential-three-eight-components}の全sourceを用いる。
磁気項は任意二次式$R_2$の全六係数である。
一般実退化因子は$D_{5,1}$。
essential\_variable\_count：3。
description：Absolutely irreducible essential-three rational cuspidal quintic; configuration E6+A4+A2。
構成は第\ref{prop:essential-three-eight-components}節、成分性は第\ref{prop:essential-three-eight-components}節、相違は第\ref{prop:essential-three-eight-components}節及びその証明を参照する。

\paragraph{族109}\label{catalogue:R4_PLANE_QUINTIC_A8_A4_23}
\path{R4_PLANE_QUINTIC_A8_A4_23}。
射影次元$23$、一般階数$4$。
一般族は命題又は構成\ref{prop:essential-three-eight-components}の全sourceを用いる。
磁気項は任意二次式$R_2$の全六係数である。
一般実退化因子は$D_{5,1}$。
essential\_variable\_count：3。
description：Absolutely irreducible essential-three rational cuspidal quintic; configuration A8+A4。
構成は第\ref{prop:essential-three-eight-components}節、成分性は第\ref{prop:essential-three-eight-components}節、相違は第\ref{prop:essential-three-eight-components}節及びその証明を参照する。

\paragraph{族110}\label{catalogue:R4_PLANE_QUINTIC_THREE_A4_23}
\path{R4_PLANE_QUINTIC_THREE_A4_23}。
射影次元$23$、一般階数$4$。
一般族は命題又は構成\ref{prop:essential-three-eight-components}の全sourceを用いる。
磁気項は任意二次式$R_2$の全六係数である。
一般実退化因子は$D_{5,1}$。
essential\_variable\_count：3。
description：Absolutely irreducible essential-three rational cuspidal quintic; configuration 3A4。
構成は第\ref{prop:essential-three-eight-components}節、成分性は第\ref{prop:essential-three-eight-components}節、相違は第\ref{prop:essential-three-eight-components}節及びその証明を参照する。

\paragraph{族111}\label{catalogue:R4_PLANE_QUINTIC_A6_THREE_A2_23}
\path{R4_PLANE_QUINTIC_A6_THREE_A2_23}。
射影次元$23$、一般階数$4$。
一般族は命題又は構成\ref{prop:essential-three-eight-components}の全sourceを用いる。
磁気項は任意二次式$R_2$の全六係数である。
一般実退化因子は$D_{5,1}$。
essential\_variable\_count：3。
description：Absolutely irreducible essential-three rational cuspidal quintic; configuration A6+3A2。
構成は第\ref{prop:essential-three-eight-components}節、成分性は第\ref{prop:essential-three-eight-components}節、相違は第\ref{prop:essential-three-eight-components}節及びその証明を参照する。

\paragraph{族112}\label{catalogue:R4_MARKED_QUARTIC_MU6_A6_CUSP_SECANT_BITANGENT_23}
\path{R4_MARKED_QUARTIC_MU6_A6_CUSP_SECANT_BITANGENT_23}。
射影次元$23$、一般階数$4$。
一般族は命題又は構成\ref{prop:marked-quartic-mu6-nine-components}の全sourceを用いる。
磁気項は任意二次式$R_2$の全六係数である。
一般実退化因子は$H_{1}+D_{4,2}$。
essential\_variable\_count：3。
description：Simple unique H plus absolutely irreducible rational essential-three plane quartic, marked case A6\_cusp\_secant\_bitangent。
構成は第\ref{prop:marked-quartic-mu6-nine-components}節、成分性は第\ref{prop:marked-quartic-mu6-nine-components}節、相違は第\ref{prop:marked-quartic-mu6-nine-components}節及びその証明を参照する。

\paragraph{族113}\label{catalogue:R4_MARKED_QUARTIC_MU6_A6_SMOOTH_FLEX_23}
\path{R4_MARKED_QUARTIC_MU6_A6_SMOOTH_FLEX_23}。
射影次元$23$、一般階数$4$。
一般族は命題又は構成\ref{prop:marked-quartic-mu6-nine-components}の全sourceを用いる。
磁気項は任意二次式$R_2$の全六係数である。
一般実退化因子は$H_{1}+D_{4,2}$。
essential\_variable\_count：3。
description：Simple unique H plus absolutely irreducible rational essential-three plane quartic, marked case A6\_smooth\_flex。
構成は第\ref{prop:marked-quartic-mu6-nine-components}節、成分性は第\ref{prop:marked-quartic-mu6-nine-components}節、相違は第\ref{prop:marked-quartic-mu6-nine-components}節及びその証明を参照する。

\paragraph{族114}\label{catalogue:R4_MARKED_QUARTIC_MU6_A4A2_CUSPS_JOIN_23}
\path{R4_MARKED_QUARTIC_MU6_A4A2_CUSPS_JOIN_23}。
射影次元$23$、一般階数$4$。
一般族は命題又は構成\ref{prop:marked-quartic-mu6-nine-components}の全sourceを用いる。
磁気項は任意二次式$R_2$の全六係数である。
一般実退化因子は$H_{1}+D_{4,2}$。
essential\_variable\_count：3。
description：Simple unique H plus absolutely irreducible rational essential-three plane quartic, marked case A4A2\_cusps\_join。
構成は第\ref{prop:marked-quartic-mu6-nine-components}節、成分性は第\ref{prop:marked-quartic-mu6-nine-components}節、相違は第\ref{prop:marked-quartic-mu6-nine-components}節及びその証明を参照する。

\paragraph{族115}\label{catalogue:R4_MARKED_QUARTIC_MU6_A4A2_A2_TANGENT_23}
\path{R4_MARKED_QUARTIC_MU6_A4A2_A2_TANGENT_23}。
射影次元$23$、一般階数$4$。
一般族は命題又は構成\ref{prop:marked-quartic-mu6-nine-components}の全sourceを用いる。
磁気項は任意二次式$R_2$の全六係数である。
一般実退化因子は$H_{1}+D_{4,2}$。
essential\_variable\_count：3。
description：Simple unique H plus absolutely irreducible rational essential-three plane quartic, marked case A4A2\_A2\_tangent。
構成は第\ref{prop:marked-quartic-mu6-nine-components}節、成分性は第\ref{prop:marked-quartic-mu6-nine-components}節、相違は第\ref{prop:marked-quartic-mu6-nine-components}節及びその証明を参照する。

\paragraph{族116}\label{catalogue:R4_MARKED_QUARTIC_MU6_A4A2_A2_SECANT_SMOOTH_TANGENT_23}
\path{R4_MARKED_QUARTIC_MU6_A4A2_A2_SECANT_SMOOTH_TANGENT_23}。
射影次元$23$、一般階数$4$。
一般族は命題又は構成\ref{prop:marked-quartic-mu6-nine-components}の全sourceを用いる。
磁気項は任意二次式$R_2$の全六係数である。
一般実退化因子は$H_{1}+D_{4,2}$。
essential\_variable\_count：3。
description：Simple unique H plus absolutely irreducible rational essential-three plane quartic, marked case A4A2\_A2\_secant\_smooth\_tangent。
構成は第\ref{prop:marked-quartic-mu6-nine-components}節、成分性は第\ref{prop:marked-quartic-mu6-nine-components}節、相違は第\ref{prop:marked-quartic-mu6-nine-components}節及びその証明を参照する。

\paragraph{族117}\label{catalogue:R4_MARKED_QUARTIC_MU6_A4A2_SMOOTH_FLEX_23}
\path{R4_MARKED_QUARTIC_MU6_A4A2_SMOOTH_FLEX_23}。
射影次元$23$、一般階数$4$。
一般族は命題又は構成\ref{prop:marked-quartic-mu6-nine-components}の全sourceを用いる。
磁気項は任意二次式$R_2$の全六係数である。
一般実退化因子は$H_{1}+D_{4,2}$。
essential\_variable\_count：3。
description：Simple unique H plus absolutely irreducible rational essential-three plane quartic, marked case A4A2\_smooth\_flex。
構成は第\ref{prop:marked-quartic-mu6-nine-components}節、成分性は第\ref{prop:marked-quartic-mu6-nine-components}節、相違は第\ref{prop:marked-quartic-mu6-nine-components}節及びその証明を参照する。

\paragraph{族118}\label{catalogue:R4_MARKED_QUARTIC_MU6_THREEA2_CUSPS_JOIN_23}
\path{R4_MARKED_QUARTIC_MU6_THREEA2_CUSPS_JOIN_23}。
射影次元$23$、一般階数$4$。
一般族は命題又は構成\ref{prop:marked-quartic-mu6-nine-components}の全sourceを用いる。
磁気項は任意二次式$R_2$の全六係数である。
一般実退化因子は$H_{1}+D_{4,2}$。
essential\_variable\_count：3。
description：Simple unique H plus absolutely irreducible rational essential-three plane quartic, marked case threeA2\_cusps\_join。
構成は第\ref{prop:marked-quartic-mu6-nine-components}節、成分性は第\ref{prop:marked-quartic-mu6-nine-components}節、相違は第\ref{prop:marked-quartic-mu6-nine-components}節及びその証明を参照する。

\paragraph{族119}\label{catalogue:R4_MARKED_QUARTIC_MU6_THREEA2_SMOOTH_BITANGENT_23}
\path{R4_MARKED_QUARTIC_MU6_THREEA2_SMOOTH_BITANGENT_23}。
射影次元$23$、一般階数$4$。
一般族は命題又は構成\ref{prop:marked-quartic-mu6-nine-components}の全sourceを用いる。
磁気項は任意二次式$R_2$の全六係数である。
一般実退化因子は$H_{1}+D_{4,2}$。
essential\_variable\_count：3。
description：Simple unique H plus absolutely irreducible rational essential-three plane quartic, marked case threeA2\_smooth\_bitangent。
構成は第\ref{prop:marked-quartic-mu6-nine-components}節、成分性は第\ref{prop:marked-quartic-mu6-nine-components}節、相違は第\ref{prop:marked-quartic-mu6-nine-components}節及びその証明を参照する。

\paragraph{族120}\label{catalogue:R4_MARKED_QUARTIC_MU6_E6_SMOOTH_BITANGENT_23}
\path{R4_MARKED_QUARTIC_MU6_E6_SMOOTH_BITANGENT_23}。
射影次元$23$、一般階数$4$。
一般族は命題又は構成\ref{prop:marked-quartic-mu6-nine-components}の全sourceを用いる。
磁気項は任意二次式$R_2$の全六係数である。
一般実退化因子は$H_{1}+D_{4,2}$。
essential\_variable\_count：3。
description：Simple unique H plus absolutely irreducible rational essential-three plane quartic, marked case E6\_smooth\_bitangent。
構成は第\ref{prop:marked-quartic-mu6-nine-components}節、成分性は第\ref{prop:marked-quartic-mu6-nine-components}節、相違は第\ref{prop:marked-quartic-mu6-nine-components}節及びその証明を参照する。

\paragraph{族121}\label{catalogue:R4_MARKED_QUARTIC_MU5_A4A1_NODE_INFLEXION_23}
\path{R4_MARKED_QUARTIC_MU5_A4A1_NODE_INFLEXION_23}。
射影次元$23$、一般階数$4$。
一般族は命題又は構成\ref{prop:marked-quartic-mu5-whole-list}の全sourceを用いる。
磁気項は任意二次式$R_2$の全六係数である。
一般実退化因子は$H_{1}+D_{4,2}$。
essential\_variable\_count：3。
description：Simple unique H plus absolutely irreducible rational essential-three plane quartic, marked case A4A1\_node\_inflexion。
構成は第\ref{prop:marked-quartic-mu5-whole-list}節、成分性は第\ref{prop:marked-quartic-mu5-four-components}節、相違は第\ref{prop:marked-quartic-mu5-four-components}節及びその証明を参照する。

\paragraph{族122}\label{catalogue:R4_MARKED_QUARTIC_MU5_A4A1_SMOOTH_HYPERFLEX_23}
\path{R4_MARKED_QUARTIC_MU5_A4A1_SMOOTH_HYPERFLEX_23}。
射影次元$23$、一般階数$4$。
一般族は命題又は構成\ref{prop:marked-quartic-mu5-whole-list}の全sourceを用いる。
磁気項は任意二次式$R_2$の全六係数である。
一般実退化因子は$H_{1}+D_{4,2}$。
essential\_variable\_count：3。
description：Simple unique H plus absolutely irreducible rational essential-three plane quartic, marked case A4A1\_smooth\_hyperflex。
構成は第\ref{prop:marked-quartic-mu5-whole-list}節、成分性は第\ref{prop:marked-quartic-mu5-four-components}節、相違は第\ref{prop:marked-quartic-mu5-four-components}節及びその証明を参照する。

\paragraph{族123}\label{catalogue:R4_MARKED_QUARTIC_MU5_D5_SMOOTH_HYPERFLEX_23}
\path{R4_MARKED_QUARTIC_MU5_D5_SMOOTH_HYPERFLEX_23}。
射影次元$23$、一般階数$4$。
一般族は命題又は構成\ref{prop:marked-quartic-mu5-whole-list}の全sourceを用いる。
磁気項は任意二次式$R_2$の全六係数である。
一般実退化因子は$H_{1}+D_{4,2}$。
essential\_variable\_count：3。
description：Simple unique H plus absolutely irreducible rational essential-three plane quartic, marked case D5\_smooth\_hyperflex。
構成は第\ref{prop:marked-quartic-mu5-whole-list}節、成分性は第\ref{prop:marked-quartic-mu5-four-components}節、相違は第\ref{prop:marked-quartic-mu5-four-components}節及びその証明を参照する。

\paragraph{族124}\label{catalogue:R4_MARKED_QUARTIC_MU5_TWOA2A1_SMOOTH_HYPERFLEX_23}
\path{R4_MARKED_QUARTIC_MU5_TWOA2A1_SMOOTH_HYPERFLEX_23}。
射影次元$23$、一般階数$4$。
一般族は命題又は構成\ref{prop:marked-quartic-mu5-whole-list}の全sourceを用いる。
磁気項は任意二次式$R_2$の全六係数である。
一般実退化因子は$H_{1}+D_{4,2}$。
essential\_variable\_count：3。
description：Simple unique H plus absolutely irreducible rational essential-three plane quartic, marked case twoA2A1\_smooth\_hyperflex。
構成は第\ref{prop:marked-quartic-mu5-whole-list}節、成分性は第\ref{prop:marked-quartic-mu5-four-components}節、相違は第\ref{prop:marked-quartic-mu5-four-components}節及びその証明を参照する。

\paragraph{族125}\label{catalogue:R4_MARKED_QUARTIC_MU6_E6_SMOOTH_FLEX_23}
\path{R4_MARKED_QUARTIC_MU6_E6_SMOOTH_FLEX_23}。
射影次元$23$、一般階数$4$。
一般族は命題又は構成\ref{prop:marked-quartic-two-normal-components}の全sourceを用いる。
磁気項は任意二次式$R_2$の全六係数である。
一般実退化因子は$H_{1}+D_{4,2}$。
essential\_variable\_count：3。
description：Simple unique H plus absolutely irreducible rational essential-three quartic, marked case E6\_smooth\_flex。
構成は第\ref{prop:marked-quartic-two-normal-components}節、成分性は第\ref{prop:marked-quartic-two-normal-components}節、相違は第\ref{prop:marked-quartic-two-normal-components}節及びその証明を参照する。

\paragraph{族126}\label{catalogue:R4_MARKED_QUARTIC_MU5_A5_SMOOTH_HYPERFLEX_23}
\path{R4_MARKED_QUARTIC_MU5_A5_SMOOTH_HYPERFLEX_23}。
射影次元$23$、一般階数$4$。
一般族は命題又は構成\ref{prop:marked-quartic-two-normal-components}の全sourceを用いる。
磁気項は任意二次式$R_2$の全六係数である。
一般実退化因子は$H_{1}+D_{4,2}$。
essential\_variable\_count：3。
description：Simple unique H plus absolutely irreducible rational essential-three quartic, marked case A5\_smooth\_hyperflex。
構成は第\ref{prop:marked-quartic-two-normal-components}節、成分性は第\ref{prop:marked-quartic-two-normal-components}節、相違は第\ref{prop:marked-quartic-two-normal-components}節及びその証明を参照する。

\paragraph{族127}\label{catalogue:R4_PLANE_DOUBLE_LINE_CUBIC_M3_TRIPLE_4_21_23}
\path{R4_PLANE_DOUBLE_LINE_CUBIC_M3_TRIPLE_4_21_23}。
射影次元$23$、一般階数$4$。
一般族は命題又は構成\ref{prop:triple-two-weak-wall-components}の全sourceを用いる。
枝 \path{tau="4/21"}。
\path{whole_magnetic_source_dimension}：6。
一般実退化因子は$D_{5,1}$。
essential\_variable\_count：5。
description：Absolutely irreducible essential-five quintic; unique ranktwo singular plane with normalorderthree and two intrinsic multiplicity-four jump lines。
構成は第\ref{prop:triple-two-weak-wall-components}節、成分性は第\ref{prop:triple-two-weak-wall-components}節、相違は第\ref{prop:triple-two-weak-wall-components}節及びその証明を参照する。

\paragraph{族128}\label{catalogue:R4_PLANE_DOUBLE_LINE_CUBIC_M3_TRIPLE_1_12_23}
\path{R4_PLANE_DOUBLE_LINE_CUBIC_M3_TRIPLE_1_12_23}。
射影次元$23$、一般階数$4$。
一般族は命題又は構成\ref{prop:triple-two-weak-wall-components}の全sourceを用いる。
枝 \path{tau="1/12"}。
\path{whole_magnetic_source_dimension}：6。
一般実退化因子は$D_{5,1}$。
essential\_variable\_count：5。
description：Absolutely irreducible essential-five quintic; unique ranktwo singular plane with normalorderthree and two intrinsic multiplicity-four jump lines。
構成は第\ref{prop:triple-two-weak-wall-components}節、成分性は第\ref{prop:triple-two-weak-wall-components}節、相違は第\ref{prop:triple-two-weak-wall-components}節及びその証明を参照する。

\paragraph{族129}\label{catalogue:R4_PLANE_DOUBLE_LINE_CUBIC_M3_MOVING_SIMPLE_2_15_23}
\path{R4_PLANE_DOUBLE_LINE_CUBIC_M3_MOVING_SIMPLE_2_15_23}。
射影次元$23$、一般階数$4$。
一般族は命題又は構成\ref{prop:moving-simple-HA-new-component}の全sourceを用いる。
\path{whole_W_dimension}：6。
一般実退化因子は$D_{5,1}$。
essential\_variable\_count：5。
description：Absolutely irreducible essentialfive quintic; unique ranktwo singularplane, normalorderthree with binaryrootpartition21.。
構成は第\ref{prop:moving-simple-HA-new-component}節、成分性は第\ref{prop:moving-simple-HA-new-component}節、相違は第\ref{prop:moving-simple-HA-new-component}節及びその証明を参照する。

\paragraph{族130}\label{catalogue:R4_MODULAR_RANK1_LONG_NONCONSTANT_25}
\path{R4_MODULAR_RANK1_LONG_NONCONSTANT_25}。
射影次元$25$、一般階数$4$。
一般族は命題又は構成\ref{prop:rankone-long-nonconstant-25}の全sourceを用いる。
\path{parameters}："g affine nonconstant in(z,w); U=u partial\_z+v partial\_w affine, d=divU, Ug=dg; K arbitrary quadratic"。
一般実退化因子は$H_{1}^{2}+H_{2}+D_{2,3}$。
essential\_variable\_count：5。
description：2H+L+Q: unique double hyperplane H tangent to smooth rank-five Q at p, with distinct simple L through p and generically not tangent to Q.。
構成は第\ref{prop:rankone-long-nonconstant-25}節、成分性は第\ref{prop:rankone-long-nonconstant-25}節、相違は第\ref{prop:rankone-long-nonconstant-25}節及びその証明を参照する。

\paragraph{族131}\label{catalogue:R4_MODULAR_RANK1_LONG_CONST_MU1_25}
\path{R4_MODULAR_RANK1_LONG_CONST_MU1_25}。
射影次元$25$、一般階数$4$。
一般族は命題又は構成\ref{prop:rankone-three-resonant-25}の全sourceを用いる。
一般実退化因子は$H_{1}^{2}+D_{3,2}$。
essential\_variable\_count：5。
description：2H+absolute irreducible cubic C; C is generically essential five.。
構成は第\ref{prop:rankone-three-resonant-25}節、成分性は第\ref{prop:rankone-three-resonant-25}節、相違は第\ref{prop:rankone-three-resonant-25}節及びその証明を参照する。

\paragraph{族132}\label{catalogue:R4_MODULAR_RANK1_LONG_CONST_MU13_25}
\path{R4_MODULAR_RANK1_LONG_CONST_MU13_25}。
射影次元$25$、一般階数$4$。
一般族は命題又は構成\ref{prop:rankone-three-resonant-25}の全sourceを用いる。
一般実退化因子は$H_{1}^{2}+D_{3,2}$。
essential\_variable\_count：5。
description：2H+absolute irreducible cubic C; C is generically essential five.。
構成は第\ref{prop:rankone-three-resonant-25}節、成分性は第\ref{prop:rankone-three-resonant-25}節、相違は第\ref{prop:rankone-three-resonant-25}節及びその証明を参照する。

\paragraph{族133}\label{catalogue:R4_MODULAR_RANK1_LONG_ZERO_SCALAR_CUBIC_25}
\path{R4_MODULAR_RANK1_LONG_ZERO_SCALAR_CUBIC_25}。
射影次元$25$、一般階数$4$。
一般族は命題又は構成\ref{prop:rankone-three-resonant-25}の全sourceを用いる。
一般実退化因子は$H_{1}^{2}+H_{2}+H_{3}+H_{4}$。
essential\_variable\_count：3。
description：2H+three distinct simple hyperplanes; simple three covectors span two, adding H gives span three.。
構成は第\ref{prop:rankone-three-resonant-25}節、成分性は第\ref{prop:rankone-three-resonant-25}節、相違は第\ref{prop:rankone-three-resonant-25}節及びその証明を参照する。

\paragraph{族134}\label{catalogue:R4_MODULAR_RANK1_GZERO_BINOMIAL_B2_25}
\path{R4_MODULAR_RANK1_GZERO_BINOMIAL_B2_25}。
射影次元$25$、一般階数$4$。
一般族は命題又は構成\ref{prop:rankone-gzero-seven-components}の全sourceを用いる。
\path{parameters}："A*B nonzero; arbitrary quadratic F with six coefficients; all actual affine q frames and centers retained"。
一般実退化因子は$H_{1}^{2}+D_{3,2}$。
essential\_variable\_count：3。
description：2H+C, C absolutely irreducible essential-three cubic; unique singular plane of C is contained in H.。
構成は第\ref{prop:rankone-gzero-seven-components}節、成分性は第\ref{prop:rankone-gzero-seven-components}節、相違は第\ref{prop:rankone-gzero-seven-components}節及びその証明を参照する。

\paragraph{族135}\label{catalogue:R4_MODULAR_RANK1_GZERO_BINOMIAL_B4_25}
\path{R4_MODULAR_RANK1_GZERO_BINOMIAL_B4_25}。
射影次元$25$、一般階数$4$。
一般族は命題又は構成\ref{prop:rankone-gzero-seven-components}の全sourceを用いる。
\path{parameters}："A*B nonzero; arbitrary quadratic F with six coefficients; all actual affine q frames and centers retained"。
一般実退化因子は$H_{1}^{2}+H_{2}+D_{2,3}$。
essential\_variable\_count：3。
description：2H+L+Q, Q rank3, restriction ranks (Q|H,Q|L)=(1,2).。
構成は第\ref{prop:rankone-gzero-seven-components}節、成分性は第\ref{prop:rankone-gzero-seven-components}節、相違は第\ref{prop:rankone-gzero-seven-components}節及びその証明を参照する。

\paragraph{族136}\label{catalogue:R4_MODULAR_RANK1_GZERO_BINOMIAL_B5_25}
\path{R4_MODULAR_RANK1_GZERO_BINOMIAL_B5_25}。
射影次元$25$、一般階数$4$。
一般族は命題又は構成\ref{prop:rankone-gzero-seven-components}の全sourceを用いる。
\path{parameters}："A*B nonzero; arbitrary quadratic F with six coefficients; all actual affine q frames and centers retained"。
一般実退化因子は$H_{1}^{2}+D_{3,2}$。
essential\_variable\_count：3。
description：2H+C, C absolutely irreducible essential-three cubic; unique singular plane of C is not contained in H.。
構成は第\ref{prop:rankone-gzero-seven-components}節、成分性は第\ref{prop:rankone-gzero-seven-components}節、相違は第\ref{prop:rankone-gzero-seven-components}節及びその証明を参照する。

\paragraph{族137}\label{catalogue:R4_MODULAR_RANK1_GZERO_BINOMIAL_B6_25}
\path{R4_MODULAR_RANK1_GZERO_BINOMIAL_B6_25}。
射影次元$25$、一般階数$4$。
一般族は命題又は構成\ref{prop:rankone-gzero-seven-components}の全sourceを用いる。
\path{parameters}："A*B nonzero; arbitrary quadratic F with six coefficients; all actual affine q frames and centers retained"。
一般実退化因子は$H_{1}^{2}+H_{2}+D_{2,3}$。
essential\_variable\_count：3。
description：2H+L+Q, Q rank3, restriction ranks (Q|H,Q|L)=(1,1).。
構成は第\ref{prop:rankone-gzero-seven-components}節、成分性は第\ref{prop:rankone-gzero-seven-components}節、相違は第\ref{prop:rankone-gzero-seven-components}節及びその証明を参照する。

\paragraph{族138}\label{catalogue:R4_MODULAR_RANK1_GZERO_B1_25}
\path{R4_MODULAR_RANK1_GZERO_B1_25}。
射影次元$25$、一般階数$4$。
一般族は命題又は構成\ref{prop:rankone-gzero-seven-components}の全sourceを用いる。
\path{parameters}："B1/B3: A*B nonzero; V1: A*(B**2-4*A*C) nonzero. Full quadratic F6 and all actual affine q frames retained."。
一般実退化因子は$H_{1}^{3}+D_{2,2}$。
essential\_variable\_count：3。
description：3H+Q rank3, Q|H rank1; unique tripleH tangent to residual conic.。
構成は第\ref{prop:rankone-gzero-seven-components}節、成分性は第\ref{prop:rankone-gzero-seven-components}節、相違は第\ref{prop:rankone-gzero-seven-components}節及びその証明を参照する。

\paragraph{族139}\label{catalogue:R4_MODULAR_RANK1_GZERO_B3_25}
\path{R4_MODULAR_RANK1_GZERO_B3_25}。
射影次元$25$、一般階数$4$。
一般族は命題又は構成\ref{prop:rankone-gzero-seven-components}の全sourceを用いる。
\path{parameters}："B1/B3: A*B nonzero; V1: A*(B**2-4*A*C) nonzero. Full quadratic F6 and all actual affine q frames retained."。
一般実退化因子は$H_{1}^{2}+H_{2}+D_{2,3}$。
essential\_variable\_count：3。
description：2H+L+Q rank3, restriction ranks(Q|H,Q|L)=(2,1).。
構成は第\ref{prop:rankone-gzero-seven-components}節、成分性は第\ref{prop:rankone-gzero-seven-components}節、相違は第\ref{prop:rankone-gzero-seven-components}節及びその証明を参照する。

\paragraph{族140}\label{catalogue:R4_MODULAR_RANK1_GZERO_V1_25}
\path{R4_MODULAR_RANK1_GZERO_V1_25}。
射影次元$25$、一般階数$4$。
一般族は命題又は構成\ref{prop:rankone-gzero-seven-components}の全sourceを用いる。
\path{parameters}："B1/B3: A*B nonzero; V1: A*(B**2-4*A*C) nonzero. Full quadratic F6 and all actual affine q frames retained."。
一般実退化因子は$H_{1}^{2}+H_{2}+H_{3}+H_{4}$。
essential\_variable\_count：3。
description：2H+three simple hyperplanes whose covectors span3; H lies in span of the two parallel affine roots.。
構成は第\ref{prop:rankone-gzero-seven-components}節、成分性は第\ref{prop:rankone-gzero-seven-components}節、相違は第\ref{prop:rankone-gzero-seven-components}節及びその証明を参照する。

\paragraph{族141}\label{catalogue:R4_MODULAR_RANK1_GZERO_M1_25}
\path{R4_MODULAR_RANK1_GZERO_M1_25}。
射影次元$25$、一般階数$4$。
一般族は命題又は構成\ref{prop:rankone-M1-M11-components}の全sourceを用いる。
\path{parameters}："t nonzero; full stated canonical coefficients, quadratic F6, and all affine q frames retained"。
一般実退化因子は$H_{1}^{4}+H_{2}$。
essential\_variable\_count：2。
description：4H+L, unique quadruple H and distinct simple L.。
構成は第\ref{prop:rankone-M1-M11-components}節、成分性は第\ref{prop:rankone-M1-M11-components}節、相違は第\ref{prop:rankone-M1-M11-components}節及びその証明を参照する。

\paragraph{族142}\label{catalogue:R4_MODULAR_RANK1_GZERO_M11_25}
\path{R4_MODULAR_RANK1_GZERO_M11_25}。
射影次元$25$、一般階数$4$。
一般族は命題又は構成\ref{prop:rankone-M1-M11-components}の全sourceを用いる。
\path{parameters}："t nonzero; full stated canonical coefficients, quadratic F6, and all affine q frames retained"。
一般実退化因子は$H_{1}^{3}+H_{2}+H_{3}$。
essential\_variable\_count：3。
description：3H+two distinct simple hyperplanes; three covectors independent.。
構成は第\ref{prop:rankone-M1-M11-components}節、成分性は第\ref{prop:rankone-M1-M11-components}節、相違は第\ref{prop:rankone-M1-M11-components}節及びその証明を参照する。

\paragraph{族143}\label{catalogue:R4_GA9_LINEAR_W_E6_CUSP_TANGENT_23}
\path{R4_GA9_LINEAR_W_E6_CUSP_TANGENT_23}。
射影次元$23$、一般階数$4$。
一般族は命題又は構成\ref{source:last-five-explicit}の全sourceを用いる。
枝 \path{source_case="Ga9linear_w"}。
一般実退化因子は$H_{1}+D_{4,2}$。
essential\_variable\_count：3。
構成は第\ref{source:last-five-explicit}節、成分性は第\ref{prop:five-last-full-components}節、相違は第\ref{prop:five-last-full-components}節及びその証明を参照する。

\paragraph{族144}\label{catalogue:R4_PLANE_TWO_LINES_CUBIC_NODE_FLEX_TANGENT_THROUGH_FLEX_23}
\path{R4_PLANE_TWO_LINES_CUBIC_NODE_FLEX_TANGENT_THROUGH_FLEX_23}。
射影次元$23$、一般階数$4$。
一般族は命題又は構成\ref{source:last-five-explicit}の全sourceを用いる。
枝 \path{source_case="node_flex_tangent_through_flex"}。
磁気項は任意二次式$R_2$の全六係数である。
一般実退化因子は$H_{1}+H_{2}+D_{3,3}$。
essential\_variable\_count：3。
構成は第\ref{source:last-five-explicit}節、成分性は第\ref{prop:five-last-full-components}節、相違は第\ref{prop:five-last-full-components}節及びその証明を参照する。

\paragraph{族145}\label{catalogue:R4_PLANE_TWO_LINES_CUBIC_CUSP_FLEX_SECANT_THROUGH_FLEX_23}
\path{R4_PLANE_TWO_LINES_CUBIC_CUSP_FLEX_SECANT_THROUGH_FLEX_23}。
射影次元$23$、一般階数$4$。
一般族は命題又は構成\ref{source:last-five-explicit}の全sourceを用いる。
枝 \path{source_case="cusp_flex_secant_through_flex"}。
磁気項は任意二次式$R_2$の全六係数である。
一般実退化因子は$H_{1}+H_{2}+D_{3,3}$。
essential\_variable\_count：3。
構成は第\ref{source:last-five-explicit}節、成分性は第\ref{prop:five-last-full-components}節、相違は第\ref{prop:five-last-full-components}節及びその証明を参照する。

\paragraph{族146}\label{catalogue:R4_PLANE_TWO_LINES_CUBIC_CUSP_CUSP_TANGENT_CUSP_SECANT_23}
\path{R4_PLANE_TWO_LINES_CUBIC_CUSP_CUSP_TANGENT_CUSP_SECANT_23}。
射影次元$23$、一般階数$4$。
一般族は命題又は構成\ref{source:last-five-explicit}の全sourceを用いる。
枝 \path{source_case="cusp_cusp_tangent_cusp_secant"}。
磁気項は任意二次式$R_2$の全六係数である。
一般実退化因子は$H_{1}+H_{2}+D_{3,3}$。
essential\_variable\_count：3。
構成は第\ref{source:last-five-explicit}節、成分性は第\ref{prop:five-last-full-components}節、相違は第\ref{prop:five-last-full-components}節及びその証明を参照する。

\paragraph{族147}\label{catalogue:R4_CONSTANT_MODULAR_DOUBLE_LINE_CUBIC_23}
\path{R4_CONSTANT_MODULAR_DOUBLE_LINE_CUBIC_23}。
射影次元$23$、一般階数$4$。
一般族は命題又は構成\ref{source:last-five-explicit}の全sourceを用いる。
枝 \path{source_case="constant_modular_highest_double"}。
一般実退化因子は$H_{1}^{2}+D_{3,2}$。
essential\_variable\_count：3。
構成は第\ref{source:last-five-explicit}節、成分性は第\ref{prop:five-last-full-components}節、相違は第\ref{prop:five-last-full-components}節及びその証明を参照する。

